\documentclass[12pt,a4paper]{article}
\usepackage{fontspec}
\usepackage{geometry}
\usepackage{enumitem}
\usepackage{fancyhdr}
\usepackage{lastpage}
\usepackage{hyperref}

% Configure IBM Plex Sans font family
\setmainfont{IBMPlexSans}[
  Path=/home/asi0/asi0-repos/bitcoin-educational-content/scripts/quizz_to_pdfs/fonts/TrueType/IBM-Plex-Sans/,
  Extension=.ttf,
  UprightFont=*-Regular,
  BoldFont=*-Bold,
  Scale=1.0
]

\geometry{margin=2.5cm}
\pagestyle{fancy}
\fancyhf{}
\rhead{Page \thepage\ of \pageref{LastPage}}
\lhead{BTC101 - Quiz (fr)}
\renewcommand{\headrulewidth}{0.4pt}

\title{Quiz: BTC101}
\author{Language: fr}
\date{\today}

\begin{document}

\maketitle
\thispagestyle{fancy}

\section*{Instructions}
Answer all questions by selecting the best option (A, B, C, or D). The answer key is provided at the end of this document.

\bigskip


\subsection*{Question 1}
Quand la guerre de la taille des blocs a-t-elle eu lieu ?

\begin{enumerate}[label=\Alph*., leftmargin=*]
\item En 2019.
\item En 2008.
\item En 2010.
\item En 2017.
\end{enumerate}

\bigskip


\subsection*{Question 2}
Quel est le but de ce cours sur Bitcoin ?

\begin{enumerate}[label=\Alph*., leftmargin=*]
\item Introduire comment créer un nouveau protocole Bitcoin.
\item Introduire comment réguler le Bitcoin.
\item Introduire comment miner du Bitcoin.
\item Introduire certains aspects monétaires généraux de Bitcoin, comment l'utiliser, et son avenir.
\end{enumerate}

\bigskip


\subsection*{Question 3}
Quelles sont les deux principales caractéristiques du Bitcoin ?

\begin{enumerate}[label=\Alph*., leftmargin=*]
\item Le Bitcoin est à la fois une monnaie physique et un portefeuille numérique.
\item Le Bitcoin est à la fois un protocole informatique et une unité monétaire.
\item Le Bitcoin est à la fois un type d'action et un type d'obligation.
\item Le Bitcoin est à la fois une monnaie centralisée et décentralisée.
\end{enumerate}

\bigskip


\subsection*{Question 4}
Quel est l'objectif principal de ce cours sur Bitcoin ?

\begin{enumerate}[label=\Alph*., leftmargin=*]
\item Les aspects monétaires de Bitcoin.
\item Les aspects de codage de Bitcoin.
\item L'impact écologique de Bitcoin.
\item Les aspects réglementaires de Bitcoin.
\end{enumerate}

\bigskip


\subsection*{Question 5}
Que signifie le fait que Bitcoin soit une monnaie neutre ?

\begin{enumerate}[label=\Alph*., leftmargin=*]
\item Elle est régulée par le gouvernement.
\item Elle n'est contrôlée par aucune entité ou institution.
\item Elle est uniquement utilisée pour de grosses transactions.
\item Elle est contrôlée par une seule entité ou institution.
\end{enumerate}

\bigskip


\subsection*{Question 6}
Qui sont les cypherpunks ?

\begin{enumerate}[label=\Alph*., leftmargin=*]
\item Un groupe de hackers qui ont volé des informations aux gouvernements.
\item Un groupe de personnes qui ont inventé le Bitcoin et d'autres cryptomonnaies.
\item Un groupe de personnes qui s'opposaient à l'utilisation de la technologie.
\item Un groupe de personnes qui ont commencé à remettre profondément en question le rôle de la vie privée et de la liberté à l'ère numérique.
\end{enumerate}

\bigskip


\subsection*{Question 7}
Qu'est-ce que le Manifeste Cypherpunk ?

\begin{enumerate}[label=\Alph*., leftmargin=*]
\item Un document décrivant les principes du Bitcoin.
\item Un livre écrit par Julian Assange sur l'avenir de la technologie.
\item Un texte écrit par Eric Hughes affirmant que la vie privée est un droit fondamental.
\item Un manifeste écrit par David Chaum sur l'importance de la monnaie numérique.
\end{enumerate}

\bigskip


\subsection*{Question 8}
Quel était le projet DigiCash ?

\begin{enumerate}[label=\Alph*., leftmargin=*]
\item Une tentative de créer de l'argent électronique anonyme par David Chaum.
\item Une tentative de numériser tout l'argent physique par la Réserve Fédérale.
\item Un projet pour créer une cryptomonnaie basée sur Bitcoin.
\item Un projet pour créer un portefeuille numérique pour Bitcoin.
\end{enumerate}

\bigskip


\subsection*{Question 9}
Qu'est-ce que la Déclaration de l'Indépendance du Cyberespace ?

\begin{enumerate}[label=\Alph*., leftmargin=*]
\item Un texte affirmant que le cyberespace devrait être contrôlé par la réglementation gouvernementale.
\item Un texte affirmant que le cyberespace est un domaine distinct de la sphère physique et ne devrait pas être soumis aux mêmes lois.
\item Une déclaration affirmant l'indépendance de Bitcoin des systèmes financiers traditionnels.
\item Une déclaration qui garantit l'indépendance du mouvement pirate de toute réglementation.
\end{enumerate}

\bigskip


\subsection*{Question 10}
Qu'est-ce que le b-money de Wei Dai ?

\begin{enumerate}[label=\Alph*., leftmargin=*]
\item Une monnaie numérique qui était largement utilisée avant Bitcoin et était contrôlée par un groupe de cypherpunks.
\item Une forme de monnaie numérique qui a été créée par la CIA et utilisée par l'armée.
\item Une idée de monnaie numérique anonyme où la détection de fraude était effectuée par une communauté d'évaluateurs plutôt que par une autorité centrale.
\item Une monnaie numérique qui a été créée par Wei Dai comme un concurrent de Bitcoin pour être plus évolutive.
\end{enumerate}

\bigskip


\subsection*{Question 11}
Que symbolise Bitcoin au-delà de sa technologie ?

\begin{enumerate}[label=\Alph*., leftmargin=*]
\item Un outil pour les criminels pour effectuer des transactions illégales.
\item Une révolution contre les systèmes financiers traditionnels et une alternative basée sur la transparence, la décentralisation et la souveraineté individuelle.
\item Une nouvelle forme de monnaie numérique qui est plus efficace que les monnaies traditionnelles car elle est plus rapide.
\item Un moyen pour les individus de gagner de l'argent par le minage.
\end{enumerate}

\bigskip


\subsection*{Question 12}
Quelles étaient les premières formes de monnaie ?

\begin{enumerate}[label=\Alph*., leftmargin=*]
\item Des biens tels que les céréales et le bétail.
\item Les actifs numériques.
\item Les billets de banque.
\item L'or et l'argent.
\end{enumerate}

\bigskip


\subsection*{Question 13}
Quelles sont les trois fonctions canoniques de la monnaie ?

\begin{enumerate}[label=\Alph*., leftmargin=*]
\item Réserve de valeur, Moyen d'échange, Unité de compte.
\item Moyen d'échange, Moyen de production, Unité de compte.
\item Réserve de valeur, Moyen de production, Unité de compte.
\item Réserve de valeur, Moyen d'échange, Moyen de production.
\end{enumerate}

\bigskip


\subsection*{Question 14}
Quelles sont les principales caractéristiques d'une monnaie efficace ?

\begin{enumerate}[label=\Alph*., leftmargin=*]
\item Fongible, Divisible, Liquide.
\item Non-fongible, Divisible, Liquide.
\item Fongible, Divisible, Périssable.
\item Fongible, Indivisible, Liquide.
\end{enumerate}

\bigskip


\subsection*{Question 15}
Quel est l'un des inconvénients de l'or en tant que forme d'argent ?

\begin{enumerate}[label=\Alph*., leftmargin=*]
\item Il est facile à produire.
\item Il n'est pas facilement divisible.
\item Il ne conserve pas sa valeur dans l'espace.
\item Il s'érode avec le temps.
\end{enumerate}

\bigskip


\subsection*{Question 16}
Pourquoi l'or est-il devenu le standard comme moyen d'échange ?

\begin{enumerate}[label=\Alph*., leftmargin=*]
\item En raison de sa rareté, de sa périssabilité et de sa divisibilité.
\item En raison de son abondance, de sa durabilité et de son indivisibilité.
\item En raison de son abondance, de sa périssabilité et de son indivisibilité.
\item En raison de sa rareté, de sa durabilité et de sa divisibilité.
\end{enumerate}

\bigskip


\subsection*{Question 17}
Quel est le rôle de l'argent comme outil de communication ?

\begin{enumerate}[label=\Alph*., leftmargin=*]
\item Il permet la communication de la valeur à travers l'espace uniquement et sans avoir besoin de faire confiance à une contrepartie.
\item Il permet la communication de la valeur à travers l'espace et le temps sans avoir besoin de faire confiance à une contrepartie.
\item Il permet la communication de la valeur à travers le temps uniquement et sans avoir besoin de faire confiance à une contrepartie.
\item Il permet la communication de la valeur à travers l'espace et le temps avec le besoin de faire confiance à une contrepartie.
\end{enumerate}

\bigskip


\subsection*{Question 18}
Quel est le principal inconvénient des monnaies fiduciaires ?

\begin{enumerate}[label=\Alph*., leftmargin=*]
\item Leur valeur est constamment appréciée par l'inflation monétaire.
\item Leur valeur est constamment érodée par l'inflation monétaire.
\item Elles ne sont pas facilement transportables.
\item Elles ne sont pas divisibles.
\end{enumerate}

\bigskip


\subsection*{Question 19}
Quel est le principal avantage du Bitcoin par rapport aux autres formes de monnaie ?

\begin{enumerate}[label=\Alph*., leftmargin=*]
\item Il offre une excellente réserve de valeur en raison de son offre strictement limitée et de sa neutralité.
\item Il est largement accepté dans le commerce par tous les commerçants du monde.
\item C'est un excellent exemple de monnaie qui subit l'inflation monétaire.
\item Il est indivisible et non-transportable.
\end{enumerate}

\bigskip


\subsection*{Question 20}
Pourquoi l'or ne peut-il pas bien performer dans un monde globalisé et numérique ?

\begin{enumerate}[label=\Alph*., leftmargin=*]
\item Parce qu'il n'est pas assez rare.
\item Parce qu'il s'érode avec le temps.
\item Parce qu'il nécessite une entité centrale pour le rendre divisible et facilement échangeable.
\item Parce qu'il n'est pas assez précieux.
\end{enumerate}

\bigskip


\subsection*{Question 21}
Pourquoi le bitcoin n'est-il pas largement accepté dans le commerce aujourd'hui ?

\begin{enumerate}[label=\Alph*., leftmargin=*]
\item Parce qu'il n'est pas divisible en petites unités.
\item Parce qu'il n'est pas liquide.
\item Parce qu'il n'est pas largement adopté par les individus.
\item Parce qu'il n'est pas transportable.
\end{enumerate}

\bigskip


\subsection*{Question 22}
Quel est le chemin d'évolution de la monnaie actuelle ?

\begin{enumerate}[label=\Alph*., leftmargin=*]
\item Pierre brute -> Pièce -> Billet de banque -> Carte bancaire -> Lightning Network
\item Pierre brute -> Pièce -> Billet de banque -> Blockchain -> Lightning Network
\item Pierre brute -> Pièce -> Carte bancaire -> Billet de banque -> Blockchain -> Lightning Network
\item Pierre brute -> Pièce -> Billet de banque -> Carte bancaire -> Blockchain -> Lightning Network
\end{enumerate}

\bigskip


\subsection*{Question 23}
Pourquoi l'or n'est-il pas un matériel portable ?

\begin{enumerate}[label=\Alph*., leftmargin=*]
\item Parce qu'il est liquide à température ambiante.
\item Parce qu'il ne peut pas être facilement façonné en formes plus petites.
\item Parce qu'il est trop précieux pour être déplacé.
\item Parce qu'il est dense et lourd, ce qui le rend encombrant à transporter en grandes quantités.
\end{enumerate}

\bigskip


\subsection*{Question 24}
Quel est le principal avantage des monnaies fiduciaires ?

\begin{enumerate}[label=\Alph*., leftmargin=*]
\item Elles sont efficaces pour conserver de la valeur dans le temps.
\item Elles sont facilement divisibles.
\item Elles ne sont pas facilement transportables.
\item Elles sont contrôlées par une partie centrale.
\end{enumerate}

\bigskip


\subsection*{Question 25}
Lequel des éléments suivants n'est pas un avantage du bitcoin en tant que monnaie ?

\begin{enumerate}[label=\Alph*., leftmargin=*]
\item Il est sans permission.
\item Il est utilisé pour payer les impôts.
\item Il est transportable.
\item Il est divisible.
\end{enumerate}

\bigskip


\subsection*{Question 26}
Comment est créée la monnaie saine ?

\begin{enumerate}[label=\Alph*., leftmargin=*]
\item La monnaie saine doit émerger naturellement du marché et du consensus volontaire.
\item La monnaie saine doit être imposée par les entreprises à travers l'utilisation de la violence.
\item La monnaie saine ne peut pas être créée.
\item La monnaie saine doit être imposée par le gouvernement par l'émission de lois monétaires.
\end{enumerate}

\bigskip


\subsection*{Question 27}
Pourquoi Bitcoin est-il considéré comme une nouvelle forme d'argent ?

\begin{enumerate}[label=\Alph*., leftmargin=*]
\item Parce que c'est une monnaie numérique sans entité centrale.
\item Parce que c'est une monnaie qui nécessite la blockchain bancaire.
\item Parce que c'est une monnaie numérique garantie par l'or.
\item Parce que la Réserve Fédérale le dit.
\end{enumerate}

\bigskip


\subsection*{Question 28}
Quel est le principal avantage de l'or en tant que monnaie ?

\begin{enumerate}[label=\Alph*., leftmargin=*]
\item Il est facilement vérifiable.
\item Il est facilement divisible ou transportable sur de longues distances.
\item Il est incroyablement durable.
\item Il n'est pas régulé par le gouvernement.
\end{enumerate}

\bigskip


\subsection*{Question 29}
Quelle est la signification de l'argent fiat ?

\begin{enumerate}[label=\Alph*., leftmargin=*]
\item L'argent fiat est une monnaie qui est soutenue par une marchandise physique, telle que l'or ou l'argent, ce qui signifie que sa valeur est directement liée à la valeur de cette marchandise.
\item L'argent fiat est une monnaie qui a de la valeur parce qu'un gouvernement la maintient et que les gens ont confiance en sa valeur, plutôt que d'être soutenue par une marchandise physique comme l'or.
\item L'argent fiat est une monnaie qui n'a pas de valeur intrinsèque et ne peut pas être utilisée pour des transactions, la rendant essentiellement sans valeur dans l'économie.
\item L'argent fiat est un type de cryptomonnaie qui fonctionne sur un réseau décentralisé, permettant des transactions de pair-à-pair sans nécessité d'une autorité centrale.
\end{enumerate}

\bigskip


\subsection*{Question 30}
Que signifie le terme fiduciaire dans le contexte des monnaies ?

\begin{enumerate}[label=\Alph*., leftmargin=*]
\item Il se réfère aux monnaies qui sont uniquement utilisées au sein d'une institution spécifique.
\item Il se réfère aux monnaies qui sont soutenues par l'or et sont émises par de petites banques commerciales.
\item Il se réfère aux monnaies qui ont de la valeur uniquement en raison de la confiance dans l'institution centrale qui les régule.
\item Il se réfère aux monnaies qui ont de la valeur uniquement pour les transactions fiduciaires.
\end{enumerate}

\bigskip


\subsection*{Question 31}
Quelles sont les entités responsables de l'émission des monnaies fiduciaires ?

\begin{enumerate}[label=\Alph*., leftmargin=*]
\item Banques commerciales.
\item Gouvernements.
\item Banques centrales.
\item Institutions privées.
\end{enumerate}

\bigskip


\subsection*{Question 32}
Quel est le nombre maximum d'unités de Bitcoin qui peuvent exister ?

\begin{enumerate}[label=\Alph*., leftmargin=*]
\item 3 billions.
\item Il n'y a pas de limite.
\item 100 millions.
\item 21 millions.
\end{enumerate}

\bigskip


\subsection*{Question 33}
Quel est l'objectif principal derrière la création du Bitcoin ?

\begin{enumerate}[label=\Alph*., leftmargin=*]
\item Éliminer l'étalon-or.
\item Rendre les transactions plus rapides et moins chères.
\item Rendre tout le monde riche.
\item Séparer l'État de l'argent.
\end{enumerate}

\bigskip


\subsection*{Question 34}
Quelle est la stratégie utilisée par les entités centrales pour dévaluer une monnaie fiduciaire ?

\begin{enumerate}[label=\Alph*., leftmargin=*]
\item Elles font fluctuer aléatoirement sa masse monétaire par rapport à la masse initiale.
\item Elles augmentent la masse monétaire de quelques pourcents chaque année.
\item Elles diminuent la masse monétaire de quelques pourcents chaque année.
\item Elles maintiennent sa masse monétaire constante par rapport à la masse initiale.
\end{enumerate}

\bigskip


\subsection*{Question 35}
Quelle est la conséquence de l'impression monétaire ?

\begin{enumerate}[label=\Alph*., leftmargin=*]
\item Cela conduit à une économie stable sans inflation ni déflation.
\item Cela conduit à la déflation, enrichissant progressivement la population.
\item Cela conduit à l'inflation, appauvrissant progressivement la population.
\item Cela conduit à une diminution de l'offre de monnaie.
\end{enumerate}

\bigskip


\subsection*{Question 36}
Quel est un inconvénient potentiel des monnaies numériques émises par les banques centrales (CBDCs) ?

\begin{enumerate}[label=\Alph*., leftmargin=*]
\item Elles pourraient entraîner une diminution de la valeur de l'or.
\item Elles pourraient entraîner une diminution de l'offre monétaire.
\item Elles pourraient entraîner une diminution de la valeur du Bitcoin.
\item Elles pourraient entraver la liberté financière des individus et faciliter les abus autoritaires.
\end{enumerate}

\bigskip


\subsection*{Question 37}
Quel est le chemin historique suivi par les institutions dans l'introduction des monnaies fiduciaires ?

\begin{enumerate}[label=\Alph*., leftmargin=*]
\item Les dirigeants introduisent une nouvelle monnaie soutenue par un panier de marchandises, qui est ensuite maintenue stable en valeur.
\item Les dirigeants introduisent une nouvelle monnaie soutenue par l'or, qui est ensuite progressivement augmentée en valeur.
\item Les dirigeants introduisent une nouvelle monnaie sans valeur intrinsèque, qui est ensuite progressivement augmentée en valeur.
\item Les dirigeants centralisent l'or et introduisent une nouvelle monnaie équivalente à l'or en valeur, qui est ensuite progressivement dévaluée.
\end{enumerate}

\bigskip


\subsection*{Question 38}
Quel est le résultat potentiel d'une dépréciation soudaine ou d'une perte de confiance dans une monnaie fiduciaire ?

\begin{enumerate}[label=\Alph*., leftmargin=*]
\item Une augmentation de la valeur de la monnaie.
\item Hyperinflation.
\item Une économie stable.
\item Déflation.
\end{enumerate}

\bigskip


\subsection*{Question 39}
Quelle est une caractéristique clé de Bitcoin qui le différencie des monnaies fiduciaires ?

\begin{enumerate}[label=\Alph*., leftmargin=*]
\item Il peut être imprimé en quantités illimitées par les banques centrales.
\item Il est soutenu par une marchandise physique comme l'or ou l'argent.
\item Il est régulé par une banque centrale.
\item Sa politique monétaire est appliquée par un consensus général.
\end{enumerate}

\bigskip


\subsection*{Question 40}
Quel est l'impact potentiel du système financier actuel sur les acteurs politiques ?

\begin{enumerate}[label=\Alph*., leftmargin=*]
\item Ils sont incités à penser à long terme.
\item Ils sont incités à effectuer des changements radicaux pour prévenir une possible implosion.
\item Ils ne sont pas incités à effectuer des changements radicaux, permettant au système de poursuivre sa trajectoire jusqu'à une possible implosion.
\item Ils ne sont pas incités à maintenir le statu quo pour assurer leur propre richesse.
\end{enumerate}

\bigskip


\subsection*{Question 41}
Qu'est-ce que l'hyperinflation ?

\begin{enumerate}[label=\Alph*., leftmargin=*]
\item Une diminution drastique de la masse monétaire.
\item Une petite diminution de la masse monétaire.
\item Une augmentation drastique de la masse monétaire.
\item Une petite augmentation de la masse monétaire.
\end{enumerate}

\bigskip


\subsection*{Question 42}
Quelle est la conséquence de l'hyperinflation ?

\begin{enumerate}[label=\Alph*., leftmargin=*]
\item Une petite augmentation de la valeur de la monnaie.
\item Une petite diminution de la valeur de la monnaie.
\item Une diminution drastique de la valeur de la monnaie.
\item Une augmentation drastique de la valeur de la monnaie.
\end{enumerate}

\bigskip


\subsection*{Question 43}
Quelle est la première phase du phénomène d'hyperinflation ?

\begin{enumerate}[label=\Alph*., leftmargin=*]
\item Perte de confiance.
\item Émergence d'une nouvelle monnaie.
\item Le cercle vicieux de l'impression monétaire.
\item Effondrement de la monnaie et augmentation des prix.
\end{enumerate}

\bigskip


\subsection*{Question 44}
Quelle est la phase finale du phénomène d'hyperinflation ?

\begin{enumerate}[label=\Alph*., leftmargin=*]
\item Le cercle vicieux de l'impression monétaire.
\item Perte de confiance.
\item Une nouvelle monnaie émerge.
\item Effondrement de la monnaie & augmentation des prix.
\end{enumerate}

\bigskip


\subsection*{Question 45}
Quel est l'impact d'une inflation de 2% sur le pouvoir d'achat sur 5 ans ?

\begin{enumerate}[label=\Alph*., leftmargin=*]
\item Vous perdez 20% de votre pouvoir d'achat.
\item Vous perdez 2% de votre pouvoir d'achat.
\item Vous perdez 5% de votre pouvoir d'achat.
\item Vous perdez 10% de votre pouvoir d'achat.
\end{enumerate}

\bigskip


\subsection*{Question 46}
Quel est l'impact d'une inflation de 20% sur le pouvoir d'achat sur 3 ans ?

\begin{enumerate}[label=\Alph*., leftmargin=*]
\item Vous perdez 50% de votre pouvoir d'achat.
\item Vous perdez 100% de votre pouvoir d'achat.
\item Vous perdez 30% de votre pouvoir d'achat.
\item Vous perdez 60% de votre pouvoir d'achat.
\end{enumerate}

\bigskip


\subsection*{Question 47}
Il existe un cas historique d'un taux d'inflation quotidien équivalent à 100 % ; où et quand cela s'est-il produit ?

\begin{enumerate}[label=\Alph*., leftmargin=*]
\item En Chine, en 1940.
\item Il s'agit en fait d'un cas irréaliste qui ne peut pas se produire.
\item Au Zimbabwe, en 2007.
\item En Hongrie, en 1945.
\end{enumerate}

\bigskip


\subsection*{Question 48}
Quelle est la deuxième phase du phénomène d'hyperinflation ?

\begin{enumerate}[label=\Alph*., leftmargin=*]
\item Effondrement de la monnaie & augmentation des prix.
\item Perte de confiance.
\item Le cercle vicieux de l'impression monétaire.
\item Émergence d'une nouvelle monnaie.
\end{enumerate}

\bigskip


\subsection*{Question 49}
Quel était approximativement le taux d'inflation en Allemagne en octobre 1923, au plus fort de l'hyperinflation ?

\begin{enumerate}[label=\Alph*., leftmargin=*]
\item 76 %
\item 5 200 %
\item 29 500 %
\item 110 %
\end{enumerate}

\bigskip


\subsection*{Question 50}
Quel était le taux d'inflation en Hongrie en juillet 1946 ?

\begin{enumerate}[label=\Alph*., leftmargin=*]
\item 207 %
\item 41 900 000 000 000 000 %
\item 100 %
\item 50 %
\end{enumerate}

\bigskip


\subsection*{Question 51}
Quel était le taux d'inflation au Zimbabwe en novembre 2008 ?

\begin{enumerate}[label=\Alph*., leftmargin=*]
\item 100 %
\item 79 600 000 000 %
\item 796 %
\item 79 600 000 000 000 000 %
\end{enumerate}

\bigskip


\subsection*{Question 52}
Que s'est-il passé avec le dollar zimbabwéen en avril 2009 ?

\begin{enumerate}[label=\Alph*., leftmargin=*]
\item Le gouvernement a dû recourir à l'utilisation de la planche à billets.
\item Le Ministre des Finances a annoncé la suspension du dollar zimbabwéen et a autorisé l'utilisation de différentes devises étrangères pour le commerce.
\item Le gouvernement a accepté de compenser les vétérans de guerre pour un montant équivalent à 450 millions de dollars américains.
\item Le dollar zimbabwéen a chuté de plus de 72%.
\end{enumerate}

\bigskip


\subsection*{Question 53}
Lequel des éléments suivants n'est pas une caractéristique du Bitcoin ?

\begin{enumerate}[label=\Alph*., leftmargin=*]
\item Il a des coûts de vérification négligeables.
\item Il a une politique monétaire fixe.
\item Il est facilement censurable.
\item Il est instantanément transportable.
\end{enumerate}

\bigskip


\subsection*{Question 54}
Quel est le processus qui vérifie les transactions sur le réseau Bitcoin ?

\begin{enumerate}[label=\Alph*., leftmargin=*]
\item Culture
\item Minage
\item Récolte
\item Fouille
\end{enumerate}

\bigskip


\subsection*{Question 55}
Quel est le nom de l'événement lors duquel la récompense de minage est divisée par deux ?

\begin{enumerate}[label=\Alph*., leftmargin=*]
\item Quartering
\item Halving
\item Tripling
\item Doubling
\end{enumerate}

\bigskip


\subsection*{Question 56}
Quel est le mécanisme qui garantit qu'un nouveau bloc est ajouté à la blockchain toutes les dix minutes ?

\begin{enumerate}[label=\Alph*., leftmargin=*]
\item Ajustement de la vitesse.
\item Ajustement du temps.
\item Ajustement de la difficulté.
\item Ajustement du bloc.
\end{enumerate}

\bigskip


\subsection*{Question 57}
Quel est le concept mathématique, basé sur la rationalité humaine, qui est utilisé dans Bitcoin ?

\begin{enumerate}[label=\Alph*., leftmargin=*]
\item Théorie des ensembles.
\item Théorie des probabilités.
\item Théorie des nombres.
\item Théorie des jeux.
\end{enumerate}

\bigskip


\subsection*{Question 58}
Quelle est la commande pour vérifier la quantité de bitcoins en circulation sur un nœud Bitcoin ?

\begin{enumerate}[label=\Alph*., leftmargin=*]
\item bitcoin-cli getblockchaininfo.
\item bitcoin-cli getnetworkinfo.
\item bitcoin-cli getwalletinfo.
\item bitcoin-cli gettxoutsetinfo.
\end{enumerate}

\bigskip


\subsection*{Question 59}
Quels principes économiques Bitcoin respecte-t-il ?

\begin{enumerate}[label=\Alph*., leftmargin=*]
\item Principes de l'économie néoclassique.
\item Principes de l'économie keynésienne.
\item Principes de l'économie marxiste.
\item Principes de l'économie autrichienne.
\end{enumerate}

\bigskip


\subsection*{Question 60}
En quelle année prévoit-on que la création de bitcoins cessera-t-elle ?

\begin{enumerate}[label=\Alph*., leftmargin=*]
\item 2050
\item 2140
\item 2200
\item 2030
\end{enumerate}

\bigskip


\subsection*{Question 61}
Quelle est la fréquence de l'ajustement de la difficulté ?

\begin{enumerate}[label=\Alph*., leftmargin=*]
\item Tous les 512 blocs.
\item Tous les 2016 blocs.
\item Tous les 1048 blocs.
\item Tous les 4032 blocs.
\end{enumerate}

\bigskip


\subsection*{Question 62}
À quelle fréquence la récompense pour le minage est-elle divisée par deux ?

\begin{enumerate}[label=\Alph*., leftmargin=*]
\item Tous les 540 000 blocs.
\item Tous les 210 000 blocs.
\item Tous les 120 000 blocs.
\item Tous les 2 100 000 blocs.
\end{enumerate}

\bigskip


\subsection*{Question 63}
Quel était le nombre estimé de bitcoins en circulation après le premier halving ?

\begin{enumerate}[label=\Alph*., leftmargin=*]
\item 5 250 000 BTC.
\item 21 000 000 BTC.
\item 10 500 000 BTC.
\item 15 750 000 BTC.
\end{enumerate}

\bigskip


\subsection*{Question 64}
Quelle est la récompense en BTC après le quatrième halving ?

\begin{enumerate}[label=\Alph*., leftmargin=*]
\item 6.25 BTC.
\item 1.5625 BTC.
\item 12.5 BTC.
\item 3.125 BTC.
\end{enumerate}

\bigskip


\subsection*{Question 65}
En quelle année aura lieu le 7ème halving ?

\begin{enumerate}[label=\Alph*., leftmargin=*]
\item 2036
\item 2048
\item 2032
\item 2028
\end{enumerate}

\bigskip


\subsection*{Question 66}
Quelles sont les trois principales fonctions d'un portefeuille Bitcoin ?

\begin{enumerate}[label=\Alph*., leftmargin=*]
\item Recevoir des bitcoins, envoyer des bitcoins, les sécuriser contre les tentatives de piratage et de vol.
\item Stocker des bitcoins, envoyer des bitcoins, les sécuriser contre les tentatives de piratage et de vol.
\item Recevoir des bitcoins, envoyer des bitcoins, vendre des bitcoins.
\item Recevoir des bitcoins, stocker des bitcoins, les sécuriser contre les tentatives de piratage et de vol.
\end{enumerate}

\bigskip


\subsection*{Question 67}
Qu'est-ce que la clé privée dans un portefeuille Bitcoin ?

\begin{enumerate}[label=\Alph*., leftmargin=*]
\item Une clé secrète pour chiffrer les transactions Bitcoin.
\item Une clé publique utilisée pour générer des adresses Bitcoin.
\item Une liste secrète de mots pour détruire un portefeuille en cas de besoin.
\item Une clé secrète qui vous permet de signer vos transactions.
\end{enumerate}

\bigskip


\subsection*{Question 68}
Où sont stockés les bitcoins ?

\begin{enumerate}[label=\Alph*., leftmargin=*]
\item Dans la blockchain Bitcoin.
\item Dans le coffre de la Réserve Fédérale.
\item Sur l'appareil où le portefeuille est installé.
\item Dans le portefeuille Bitcoin.
\end{enumerate}

\bigskip


\subsection*{Question 69}
Qu'est-ce que la phrase mnémonique dans un portefeuille Bitcoin ?

\begin{enumerate}[label=\Alph*., leftmargin=*]
\item Une liste de 12 ou 24 mots qui vous permet d'accéder à vos bitcoins.
\item Une phrase de récupération secrète qui détruit le portefeuille Bitcoin en cas de besoin.
\item Un code pour crypter les transactions Bitcoin.
\item Un mot de passe choisi pour accéder au portefeuille Bitcoin.
\end{enumerate}

\bigskip


\subsection*{Question 70}
Quel est le rôle de la clé publique dans un portefeuille Bitcoin ?

\begin{enumerate}[label=\Alph*., leftmargin=*]
\item Elle est utilisée pour chiffrer les transactions Bitcoin.
\item Elle est utilisée pour sauvegarder le portefeuille Bitcoin.
\item Elle est utilisée pour accéder au portefeuille Bitcoin.
\item Elle est utilisée pour générer des adresses Bitcoin et ainsi recevoir de l'argent.
\end{enumerate}

\bigskip


\subsection*{Question 71}
Quel est le risque de partager la clé publique d'un portefeuille Bitcoin ?

\begin{enumerate}[label=\Alph*., leftmargin=*]
\item Perdre la confidentialité.
\item Détruire le portefeuille Bitcoin.
\item Perdre les Bitcoins.
\item Pirater le portefeuille Bitcoin.
\end{enumerate}

\bigskip


\subsection*{Question 72}
Laquelle de ces options est fausse concernant la perte de l'appareil sur lequel se trouve votre portefeuille Bitcoin ?

\begin{enumerate}[label=\Alph*., leftmargin=*]
\item Vous pouvez recréer votre portefeuille et dépenser vos bitcoins en réinitialisant votre mot de passe.
\item Vous pouvez recréer votre portefeuille et dépenser vos bitcoins en utilisant la phrase mnémonique.
\item Vous pouvez vérifier avec vos clés publiques que vos bitcoins n'ont pas été dépensés.
\item Vous perdez tous vos bitcoins si vous n'avez pas la phrase mnémonique.
\end{enumerate}

\bigskip


\subsection*{Question 73}
Comment pouvez-vous renforcer la sécurité par défaut de votre portefeuille Bitcoin ?

\begin{enumerate}[label=\Alph*., leftmargin=*]
\item En ajoutant une seconde phrase mnémonique à votre portefeuille Bitcoin.
\item En ajoutant une seconde clé privée à votre portefeuille Bitcoin.
\item En ajoutant une phrase secrète à votre portefeuille Bitcoin.
\item En ajoutant une seconde clé publique à votre portefeuille Bitcoin.
\end{enumerate}

\bigskip


\subsection*{Question 74}
Quelle est la probabilité que quelqu'un devine accidentellement votre liste de 12 ou 24 mots dans un portefeuille Bitcoin ?

\begin{enumerate}[label=\Alph*., leftmargin=*]
\item Modérée.
\item Élevée.
\item Forte.
\item Extrêmement faible.
\end{enumerate}

\bigskip


\subsection*{Question 75}
Quelle est la question centrale à se poser lors du choix d'un portefeuille Bitcoin, en termes de garde personnelle ?

\begin{enumerate}[label=\Alph*., leftmargin=*]
\item Êtes-vous le propriétaire des fonds ou laissez-vous le contrôle de votre argent à un tiers ?
\item Le portefeuille permet-il de recevoir et d'envoyer des bitcoins ?
\item Puis-je partager mon portefeuille avec une autre personne ?
\item Le portefeuille est-il sécurisé par une phrase secrète ?
\end{enumerate}

\bigskip


\subsection*{Question 76}
Quelle est la technologie utilisée pour dépenser des fonds de manière sécurisée en utilisant un portefeuille Bitcoin ?

\begin{enumerate}[label=\Alph*., leftmargin=*]
\item Cryptographie RSA
\item Cryptographie sur courbes elliptiques
\item Cryptographie quantique
\item Cryptographie symétrique
\end{enumerate}

\bigskip


\subsection*{Question 77}
Quelle est l'analogie utilisée dans ce cours pour expliquer comment vous pouvez recevoir des bitcoins sans permettre à l'expéditeur de voler vos fonds ?

\begin{enumerate}[label=\Alph*., leftmargin=*]
\item Une boîte aux lettres
\item Un coffre-fort
\item Une maison
\item Un compte bancaire
\end{enumerate}

\bigskip


\subsection*{Question 78}
Quelle devrait être la principale préoccupation lorsque vous possédez des bitcoins ?

\begin{enumerate}[label=\Alph*., leftmargin=*]
\item La valeur du Bitcoin en termes de dollars.
\item La permission de dépenser vos fonds.
\item La légalité de posséder des bitcoins.
\item La sécurité de vos fonds.
\end{enumerate}

\bigskip


\subsection*{Question 79}
Qu'est-ce qu'un portefeuille chaud ?

\begin{enumerate}[label=\Alph*., leftmargin=*]
\item Un portefeuille qui est recommandé pour stocker de grandes quantités de bitcoins.
\item Un portefeuille qui est physiquement chaud au toucher.
\item Un portefeuille qui est populaire parmi les utilisateurs de Bitcoin.
\item Un portefeuille Bitcoin sur votre téléphone ou ordinateur connecté à internet.
\end{enumerate}

\bigskip


\subsection*{Question 80}
Qu'est-ce qu'un portefeuille froid ?

\begin{enumerate}[label=\Alph*., leftmargin=*]
\item Un portefeuille dans un dispositif physique qui stocke vos clés et qui n'est pas connecté à internet.
\item Un portefeuille qui est recommandé pour stocker de petites quantités de bitcoins.
\item Un portefeuille qui est physiquement froid au toucher.
\item Un portefeuille qui est impopulaire parmi les utilisateurs de Bitcoin.
\end{enumerate}

\bigskip


\subsection*{Question 81}
Qu'est-ce qu'un portefeuille multisig ?

\begin{enumerate}[label=\Alph*., leftmargin=*]
\item Un portefeuille qui nécessite une signature physique pour effectuer une transaction.
\item Un portefeuille qui ne nécessite aucune signature pour effectuer une transaction.
\item Un ensemble de portefeuilles qui nécessite plusieurs signatures pour effectuer une transaction.
\item Un portefeuille qui nécessite une seule signature pour effectuer une transaction.
\end{enumerate}

\bigskip


\subsection*{Question 82}
Quel est le principal risque d'utiliser un service de garde pour stocker vos bitcoins ?

\begin{enumerate}[label=\Alph*., leftmargin=*]
\item Le tiers de confiance peut augmenter la valeur de vos bitcoins à tout moment.
\item Le tiers de confiance peut diminuer la valeur de vos bitcoins à tout moment.
\item Le tiers de confiance ne peut pas dupliquer vos bitcoins à tout moment.
\item Le tiers de confiance peut restreindre votre accès à vos fonds à tout moment.
\end{enumerate}

\bigskip


\subsection*{Question 83}
Quel est le risque de compliquer excessivement la sécurité et l'accessibilité de vos bitcoins ?

\begin{enumerate}[label=\Alph*., leftmargin=*]
\item Vous pouvez vous nuire si vous manipulez mal les portefeuilles de sauvegarde.
\item Vous pouvez vous nuire si vous manipulez mal les portefeuilles numériques.
\item Vous pouvez vous nuire si vous manipulez mal les portefeuilles multisig.
\item Vous pouvez vous nuire si vous manipulez mal les portefeuilles physiques.
\end{enumerate}

\bigskip


\subsection*{Question 84}
Quel est l'usage recommandé d'un portefeuille mobile ?

\begin{enumerate}[label=\Alph*., leftmargin=*]
\item Stocker les dépenses quotidiennes.
\item Stocker les économies à moyen terme.
\item Stocker de grandes sommes.
\item Stocker les économies à long terme.
\end{enumerate}

\bigskip


\subsection*{Question 85}
Quel est l'usage recommandé d'un portefeuille physique hors ligne, ou cold wallet ?

\begin{enumerate}[label=\Alph*., leftmargin=*]
\item Stocker de petites quantités.
\item Stocker des économies à court terme.
\item Stocker les dépenses quotidiennes.
\item Stocker de grandes quantités.
\end{enumerate}

\bigskip


\subsection*{Question 86}
Quel est l'usage recommandé d'un portefeuille multisig ?

\begin{enumerate}[label=\Alph*., leftmargin=*]
\item Stocker les dépenses quotidiennes pour un usage individuel.
\item Stocker des quantités moyennes pour un usage à but non lucratif.
\item Stocker de petites quantités pour un usage personnel.
\item Stocker de grandes quantités pour un usage corporatif.
\end{enumerate}

\bigskip


\subsection*{Question 87}
De quoi avez-vous besoin pour effectuer une sauvegarde lorsque vous utilisez un portefeuille avec une phrase secrète ?

\begin{enumerate}[label=\Alph*., leftmargin=*]
\item Vous n'avez besoin de sauvegarder ni la liste de 12 ou 24 mots, ni votre phrase secrète.
\item Vous devez seulement sauvegarder la liste de 12 ou 24 mots.
\item Vous devez sauvegarder à la fois la liste de 12 ou 24 mots et votre phrase secrète.
\item Vous devez seulement sauvegarder votre phrase secrète.
\end{enumerate}

\bigskip


\subsection*{Question 88}
Quel est le risque d'utiliser un portefeuille multisig ?

\begin{enumerate}[label=\Alph*., leftmargin=*]
\item Partager vos clés privées avec d'autres parties.
\item Avoir accès à suffisamment de clés privées individuelles ou de composants de sauvegarde.
\item Partager vos clés publiques avec d'autres parties.
\item Ne pas avoir accès à suffisamment de clés privées individuelles ou de composants de sauvegarde.
\end{enumerate}

\bigskip


\subsection*{Question 89}
Que signifie le terme hodl ?

\begin{enumerate}[label=\Alph*., leftmargin=*]
\item Hodl est un terme d'argot dans la communauté des cryptomonnaies qui signifie conserver vos investissements au lieu de les vendre, surtout pendant la volatilité du marché.
\item Hodl est une stratégie de trading journalier qui implique d'acheter et de vendre fréquemment.
\item Hodl fait référence à un type de cryptomonnaie qui est utilisé uniquement pour le trading.
\item Hodl est un terme qui signifie vendre toute votre cryptomonnaie à un prix élevé.
\end{enumerate}

\bigskip


\subsection*{Question 90}
Par quoi la clé privée est-elle souvent représentée dans un portefeuille Bitcoin ?

\begin{enumerate}[label=\Alph*., leftmargin=*]
\item Une liste de 12 ou 48 mots.
\item Une liste de 12 ou 24 mots.
\item Une liste de 32 mots.
\item Une liste de 64 mots.
\end{enumerate}

\bigskip


\subsection*{Question 91}
Que se passe-t-il si votre clé privée est révélée à un tiers ?

\begin{enumerate}[label=\Alph*., leftmargin=*]
\item Le tiers peut dépenser vos bitcoins.
\item Le tiers peut uniquement consulter votre historique de transactions.
\item Le tiers peut changer votre clé privée.
\item Le tiers peut accéder à vos informations personnelles.
\end{enumerate}

\bigskip


\subsection*{Question 92}
Quelle est une alternative recommandée à l'utilisation du papier pour stocker votre phrase mnémonique ?

\begin{enumerate}[label=\Alph*., leftmargin=*]
\item La stocker dans un fichier numérique.
\item L'écrire sur un tableau blanc.
\item La sauvegarder comme un contact dans votre téléphone.
\item La graver sur une plaque métallique.
\end{enumerate}

\bigskip


\subsection*{Question 93}
Que devez-vous faire une fois les mots de votre phrase mnémonique écrits ?

\begin{enumerate}[label=\Alph*., leftmargin=*]
\item Partagez-la avec un ami de confiance.
\item Stockez-la dans votre téléphone.
\item Envoyez-la par courriel à vous-même.
\item Faites-en une seconde copie et stockez-la dans un endroit séparé.
\end{enumerate}

\bigskip


\subsection*{Question 94}
Que signifie l'absence d'une phrase mnémonique dans un portefeuille ?

\begin{enumerate}[label=\Alph*., leftmargin=*]
\item Cela indique que le portefeuille est entièrement sécurisé et ne nécessite aucune sauvegarde.
\item Cela suggère que le portefeuille est automatiquement sauvegardé par votre appareil.
\item Cela signifie que le portefeuille peut contenir de nombreux types d'obligations différents.
\item Cela indique que le portefeuille n'est pas sécurisé pour le stockage à long terme de Bitcoin.
\end{enumerate}

\bigskip


\subsection*{Question 95}
Quelle est la méthode standard pour restaurer votre portefeuille ?

\begin{enumerate}[label=\Alph*., leftmargin=*]
\item Saisir votre phrase mnémonique dans n'importe quel logiciel de portefeuille.
\item Écrire votre phrase mnémonique sur un tableau blanc.
\item Sauvegarder votre phrase mnémonique dans un fichier numérique.
\item Partager votre phrase mnémonique avec un ami de confiance.
\end{enumerate}

\bigskip


\subsection*{Question 96}
Quelle est une méthode pour sécuriser vos bitcoins sur le long terme ?

\begin{enumerate}[label=\Alph*., leftmargin=*]
\item Partager votre phrase mnémonique avec un collègue.
\item Graver votre phrase mnémonique sur un matériau résistant comme l'acier.
\item Les stocker dans un portefeuille en ligne.
\item Les conserver sur une plateforme d'échange.
\end{enumerate}

\bigskip


\subsection*{Question 97}
Quel est le but de créer un plan de succession pour vos bitcoins ?

\begin{enumerate}[label=\Alph*., leftmargin=*]
\item Pour éviter les impôts sur vos bitcoins.
\item Pour empêcher vos bitcoins d'être volés.
\item Pour garantir que vos bitcoins seront correctement gérés après votre décès.
\item Pour diminuer la valeur de vos bitcoins.
\end{enumerate}

\bigskip


\subsection*{Question 98}
Qu'est-ce qui n'est pas une phrase mnémonique dans le contexte du Bitcoin ?

\begin{enumerate}[label=\Alph*., leftmargin=*]
\item Une liste de 24 mots issus d'un dictionnaire spécifique.
\item Un mot de passe pour votre portefeuille Bitcoin.
\item Le secret pour accéder à vos fonds en bitcoin.
\item Une liste de 12 mots issus d'un dictionnaire spécifique.
\end{enumerate}

\bigskip


\subsection*{Question 99}
Pourquoi la confidentialité est-elle importante dans le contexte du Bitcoin ?

\begin{enumerate}[label=\Alph*., leftmargin=*]
\item Pour diminuer la valeur de vos bitcoins.
\item Pour payer des impôts sur vos bitcoins.
\item Pour empêcher vos bitcoins d'être volés.
\item Pour minimiser le risque que vos fonds soient tracés.
\end{enumerate}

\bigskip


\subsection*{Question 100}
Pourquoi est-il conseillé d'éviter de laisser vos bitcoins sur des plateformes d'échange ?

\begin{enumerate}[label=\Alph*., leftmargin=*]
\item Ils peuvent être susceptibles aux attaques de hackers.
\item Les plateformes d'échange imposent des frais aléatoires.
\item Les plateformes d'échange ne sont pas des institutions régulées.
\item Le bitcoin perd de sa valeur sur les plateformes d'échange.
\end{enumerate}

\bigskip


\subsection*{Question 101}
Quel est le rôle d'un notaire dans le contexte de l'héritage de Bitcoin ?

\begin{enumerate}[label=\Alph*., leftmargin=*]
\item Ils assurent la conformité fiscale.
\item Ils gèrent vos bitcoins après votre décès.
\item Ils modifient votre plan d'héritage.
\item Ils sécurisent vos bitcoins dans un coffre-fort.
\end{enumerate}

\bigskip


\subsection*{Question 102}
Quel est le but d'un portefeuille Bitcoin ?

\begin{enumerate}[label=\Alph*., leftmargin=*]
\item Utiliser les ressources de l'appareil pour miner des bitcoins.
\item Stocker des bitcoins sur un appareil.
\item Stocker les clés privées afin de réaliser des transactions en bitcoins.
\item Acheter des bitcoins via la clé publique.
\end{enumerate}

\bigskip


\subsection*{Question 103}
Quelle est l'importance de la responsabilité individuelle dans Bitcoin ?

\begin{enumerate}[label=\Alph*., leftmargin=*]
\item Elle est essentielle pour acheter des bitcoins.
\item Elle est cruciale parce que vous êtes responsable des sauvegardes.
\item Elle est essentielle parce que vous devez miner des bitcoins.
\item Elle est cruciale parce que vous êtes responsable du thème du portefeuille.
\end{enumerate}

\bigskip


\subsection*{Question 104}
Qu'est-ce qu'un Blockmit dans le contexte de la sécurité Bitcoin ?

\begin{enumerate}[label=\Alph*., leftmargin=*]
\item Un type de portefeuille Bitcoin qui nécessite plusieurs signatures pour dépenser les UTXOs.
\item Une solution peu coûteuse pour graver votre phrase mnémonique sur un morceau d'acier.
\item Une solution peu coûteuse pour miner des bitcoins à domicile.
\item Un type de portefeuille Bitcoin pour signer plusieurs transactions à la fois.
\end{enumerate}

\bigskip


\subsection*{Question 105}
Que signifie KYC ?

\begin{enumerate}[label=\Alph*., leftmargin=*]
\item Gardez Votre Argent.
\item Veuillez Vous Conformer Gentiment.
\item Vos Clés Vos Codes.
\item Connaître Votre Client.
\end{enumerate}

\bigskip


\subsection*{Question 106}
Pourquoi devriez-vous acheter des bitcoins non-KYC ?

\begin{enumerate}[label=\Alph*., leftmargin=*]
\item Pour garantir votre vie privée.
\item Pour vendre à bas prix.
\item Pour garantir l'enregistrement de tous les utilisateurs de Bitcoin.
\item Pour acheter à bas prix.
\end{enumerate}

\bigskip


\subsection*{Question 107}
Qu'est-ce qui empêche tout le monde d'utiliser le même portefeuille Bitcoin ?

\begin{enumerate}[label=\Alph*., leftmargin=*]
\item Les utilisateurs ont des besoins différents selon les cas d'utilisation.
\item Il est illégal d'utiliser un portefeuille Bitcoin.
\item Il est coûteux d'utiliser un portefeuille Bitcoin.
\item Les gens ne seront jamais d'accord sur le nom d'un tel portefeuille.
\end{enumerate}

\bigskip


\subsection*{Question 108}
Quand le Bitcoin a-t-il été présenté au monde pour la première fois ?

\begin{enumerate}[label=\Alph*., leftmargin=*]
\item 22 novembre 2009
\item 31 octobre 2008
\item 12 décembre 2010
\item 3 janvier 2009
\end{enumerate}

\bigskip


\subsection*{Question 109}
Qui était la première personne à rejoindre le réseau Bitcoin après Satoshi Nakamoto ?

\begin{enumerate}[label=\Alph*., leftmargin=*]
\item Hal Finney
\item Phil Champagne
\item Laszlo Hanyecz
\item Rogzy
\end{enumerate}

\bigskip


\subsection*{Question 110}
Quelle est la première transaction Bitcoin enregistrée ?

\begin{enumerate}[label=\Alph*., leftmargin=*]
\item Une transaction de 1 BTC entre Satoshi Nakamoto et Hal Finney.
\item Une transaction de 10 BTC entre Satoshi Nakamoto et Hal Finney.
\item Une transaction de 10 000 BTC pour 2 pizzas.
\item Une transaction de 10 000 BTC entre Laszlo Hanyecz et Hal Finney.
\end{enumerate}

\bigskip


\subsection*{Question 111}
Quand Satoshi Nakamoto a-t-il disparu ?

\begin{enumerate}[label=\Alph*., leftmargin=*]
\item En 2023
\item En 2009
\item En 2010
\item En 2008
\end{enumerate}

\bigskip


\subsection*{Question 112}
Quel est le message contenu dans le premier bloc de la blockchain Bitcoin ?

\begin{enumerate}[label=\Alph*., leftmargin=*]
\item 03/jan/2009 Le Chancelier au bord d'un second sauvetage des banques.
\item Les gouvernements sont doués pour couper la tête de réseaux centralisés comme Napster, mais les réseaux purement P2P comme Gnutella et Tor semblent tenir le coup.
\item Aucune des réponses ci-dessus.
\item Nous pouvons remporter une bataille majeure dans la course aux armements et gagner un nouveau territoire de liberté pour plusieurs années.
\end{enumerate}

\bigskip


\subsection*{Question 113}
Qu'est-ce que BitcoinTalk ?

\begin{enumerate}[label=\Alph*., leftmargin=*]
\item Un portefeuille Bitcoin.
\item Un logiciel de minage de Bitcoin.
\item Une conférence sur le Bitcoin.
\item Un lieu de discussion pour les utilisateurs de Bitcoin.
\end{enumerate}

\bigskip


\subsection*{Question 114}
Quand est-ce que le Bitcoin a été utilisé pour la première fois pour acheter des biens physiques ?

\begin{enumerate}[label=\Alph*., leftmargin=*]
\item Quand Hal Finney a acheté une voiture pour 100 000 BTC.
\item Aucune des autres options.
\item Quand Laszlo Hanyecz a acheté 2 pizzas pour 10 000 BTC.
\item Quand Satoshi Nakamoto a acheté une maison pour 1 000 000 000 BTC.
\end{enumerate}

\bigskip


\subsection*{Question 115}
Depuis sa création, quel pourcentage du temps le réseau Bitcoin a-t-il été fonctionnel ?

\begin{enumerate}[label=\Alph*., leftmargin=*]
\item 99,988 %
\item 80,17 %
\item 90,009 %
\item 100 %
\end{enumerate}

\bigskip


\subsection*{Question 116}
Quelle a été la disponibilité en ligne de Bitcoin depuis 2009 ?

\begin{enumerate}[label=\Alph*., leftmargin=*]
\item 99,988 %
\item 100 %
\item 90,009 %
\item 80,17 %
\end{enumerate}

\bigskip


\subsection*{Question 117}
Quel est le numéro du bloc de la première transaction Bitcoin ?

\begin{enumerate}[label=\Alph*., leftmargin=*]
\item Bloc 170
\item Bloc 1
\item Bloc 17
\item Bloc 1700
\end{enumerate}

\bigskip


\subsection*{Question 118}
Quelle est la version du premier logiciel Bitcoin publié ?

\begin{enumerate}[label=\Alph*., leftmargin=*]
\item Bitcoin-1.0.0
\item Bitcoin-0.1.1
\item Bitcoin-0.1.0
\item Bitcoin-0.0.1
\end{enumerate}

\bigskip


\subsection*{Question 119}
Quelle est la date du dernier message privé connu de Satoshi Nakamoto par email ?

\begin{enumerate}[label=\Alph*., leftmargin=*]
\item 12 décembre 2010
\item 31 octobre 2008
\item 8 janvier 2009
\item 23 avril 2011
\end{enumerate}

\bigskip


\subsection*{Question 120}
Qu'est-ce qu'une transaction Bitcoin ?

\begin{enumerate}[label=\Alph*., leftmargin=*]
\item Un transfert de propriété de bitcoins, utilisant une adresse email.
\item Un transfert de propriété de bitcoins, utilisant une adresse physique.
\item Un transfert de propriété de bitcoins, utilisant la cryptographie.
\item Un échange physique de bitcoins, utilisant un notaire.
\end{enumerate}

\bigskip


\subsection*{Question 121}
Quel est le but des frais de transaction dans les transactions Bitcoin ?

\begin{enumerate}[label=\Alph*., leftmargin=*]
\item Ils sont des frais facturés par le fournisseur de portefeuille.
\item Ils sont des frais facturés par le réseau Bitcoin pour l'utilisation de ses services.
\item Ils constituent une incitation pour les mineurs à inclure la transaction dans le bloc suivant.
\item Ils représentent une taxe imposée par le gouvernement.
\end{enumerate}

\bigskip


\subsection*{Question 122}
Lorsque vous créez un plan de succession, qu'est-ce qui n'est pas nécessaire d'inclure ?

\begin{enumerate}[label=\Alph*., leftmargin=*]
\item Une lettre manuscrite listant vos actifs.
\item Une description de comment accéder à vos fonds.
\item Les informations de contact de personnes de confiance.
\item Le nom du portefeuille Bitcoin que vous utilisez habituellement.
\end{enumerate}

\bigskip


\subsection*{Question 123}
Qu'est-ce que la blockchain Bitcoin ?

\begin{enumerate}[label=\Alph*., leftmargin=*]
\item Un registre public et immuable de toutes les transactions Bitcoin.
\item Un registre privé et mutable de toutes les transactions Bitcoin.
\item Un registre public et mutable de toutes les transactions Bitcoin.
\item Un registre privé et immuable de toutes les transactions Bitcoin.
\end{enumerate}

\bigskip


\subsection*{Question 124}
Pourquoi le nombre de confirmations dans une transaction Bitcoin est-il important ?

\begin{enumerate}[label=\Alph*., leftmargin=*]
\item Plus une transaction a de confirmations, plus elle devient immuable.
\item Plus une transaction a de confirmations, moins elle est sécurisée.
\item Plus une transaction a de confirmations, plus elle est traitée rapidement.
\item Plus une transaction a de confirmations, plus elle devient précieuse.
\end{enumerate}

\bigskip


\subsection*{Question 125}
Quel est le rôle d'un nœud Bitcoin dans le réseau Bitcoin ?

\begin{enumerate}[label=\Alph*., leftmargin=*]
\item Il gère le réseau Bitcoin.
\item Il crée de nouveaux Bitcoins et les place dans les blocs.
\item Il régule le prix du Bitcoin.
\item Il valide les nouveaux blocs et les transactions incluses.
\end{enumerate}

\bigskip


\subsection*{Question 126}
Quelle est l'importance de la répartition géographique des nœuds Bitcoin ?

\begin{enumerate}[label=\Alph*., leftmargin=*]
\item Cela affecte la vitesse des transactions Bitcoin.
\item Cela rend pratiquement impossible de détruire toutes les copies de la blockchain.
\item Cela détermine le prix du Bitcoin dans différentes régions.
\item Cela détermine le nombre de bitcoins qui peuvent être minés dans une région.
\end{enumerate}

\bigskip


\subsection*{Question 127}
Comment les mineurs valident-ils les transactions et les incluent-ils dans un bloc ?

\begin{enumerate}[label=\Alph*., leftmargin=*]
\item Par le processus de Preuve de destruction (Proof of burn).
\item Par le processus de Preuve de travail (Proof of work).
\item Par le processus de Preuve d'autorité (Proof of authority).
\item Par le processus de Preuve d'enjeu (Proof of stake).
\end{enumerate}

\bigskip


\subsection*{Question 128}
Qu'est-ce qu'un hash valide dans le contexte du minage de blocs dans le réseau Bitcoin ?

\begin{enumerate}[label=\Alph*., leftmargin=*]
\item Un identifiant unique pour le mineur qui a miné le bloc.
\item Un code unique utilisé pour accéder au bloc.
\item Une empreinte digitale unique associée au bloc, composée de 256 caractères.
\item Un mot de passe unique requis pour miner le bloc.
\end{enumerate}

\bigskip


\subsection*{Question 129}
Quelle est la fonction de l'ajustement de difficulté dans le réseau Bitcoin ?

\begin{enumerate}[label=\Alph*., leftmargin=*]
\item Il garantit qu'un bloc est ajouté à la blockchain environ toutes les dix minutes.
\item Il détermine la taille d'un bloc.
\item Il détermine le prix du Bitcoin.
\item Il régule le nombre de bitcoins qui peuvent être minés dans la blockchain.
\end{enumerate}

\bigskip


\subsection*{Question 130}
Comment le réseau Bitcoin s'assure-t-il qu'il est trop coûteux d'agir de manière malveillante au sein du protocole Bitcoin ?

\begin{enumerate}[label=\Alph*., leftmargin=*]
\item Par la menace d'action légale.
\item Par des règles de protocole et des incitations économiques.
\item Par l'utilisation de mesures de sécurité avancées.
\item Par une régulation stricte et une supervision.
\end{enumerate}

\bigskip


\subsection*{Question 131}
Comment les utilisateurs transfèrent-ils la propriété de leur argent dans le réseau Bitcoin ?

\begin{enumerate}[label=\Alph*., leftmargin=*]
\item En fournissant leur adresse Bitcoin au destinataire.
\item En signant numériquement les transactions avec leurs clés privées.
\item En envoyant un email au destinataire.
\item En remettant physiquement leurs bitcoins.
\end{enumerate}

\bigskip


\subsection*{Question 132}
Quelle est la fonction principale des nœuds dans le réseau Bitcoin ?

\begin{enumerate}[label=\Alph*., leftmargin=*]
\item Maintenir une copie de la blockchain Bitcoin, valider les transactions, transmettre les informations aux autres nœuds et faire respecter les règles du protocole Bitcoin.
\item Transmettre uniquement les informations aux autres nœuds.
\item Maintenir une copie de la blockchain Bitcoin, valider les transactions, transmettre les informations au gouvernement et modifier les règles du protocole Bitcoin.
\item Maintenir uniquement une copie de la blockchain Bitcoin.
\end{enumerate}

\bigskip


\subsection*{Question 133}
Quel est le nombre approximatif de nœuds répartis à travers le monde en septembre 2023 ?

\begin{enumerate}[label=\Alph*., leftmargin=*]
\item 45 000 nœuds
\item 100 000 nœuds
\item 10 000 nœuds
\item 500 000 nœuds
\end{enumerate}

\bigskip


\subsection*{Question 134}
Quelle est la taille actuelle de la blockchain Bitcoin ?

\begin{enumerate}[label=\Alph*., leftmargin=*]
\item Environ 110Go
\item Environ 2To
\item Environ 660Go
\item Environ 1To
\end{enumerate}

\bigskip


\subsection*{Question 135}
Quel est le coût de fonctionnement d'un nœud Bitcoin sur un Raspberry Pi 4 en termes de consommation électrique par an (en 2023) ?

\begin{enumerate}[label=\Alph*., leftmargin=*]
\item Un peu plus de 50€ par an.
\item Un peu moins de 100€ par an.
\item Environ 200€ par an.
\item Un peu moins de 10€ par an.
\end{enumerate}

\bigskip


\subsection*{Question 136}
Qu'est-ce qu'un nœud élagué dans le contexte du Bitcoin ?

\begin{enumerate}[label=\Alph*., leftmargin=*]
\item Un nœud qui transmet uniquement des informations aux autres nœuds.
\item Un nœud qui ne conserve que les derniers N blocs en mémoire.
\item Un nœud qui valide uniquement les transactions.
\item Un nœud qui applique uniquement les règles du protocole Bitcoin.
\end{enumerate}

\bigskip


\subsection*{Question 137}
Que se passerait-il si un utilisateur de nœud décidait de suivre des règles différentes qui contredisent les règles de consensus existantes ?

\begin{enumerate}[label=\Alph*., leftmargin=*]
\item Le nœud serait récompensé par le réseau Bitcoin.
\item Le nœud deviendrait une partie d'un réseau Bitcoin différent.
\item Le nœud serait pénalisé par le réseau Bitcoin.
\item Le nœud ne ferait plus partie du réseau Bitcoin.
\end{enumerate}

\bigskip


\subsection*{Question 138}
Quelle est l'exigence de bande passante pour exécuter un nœud Bitcoin, en considérant 1 bloc de 1MB toutes les 10 minutes ?

\begin{enumerate}[label=\Alph*., leftmargin=*]
\item Environ 5GB par mois.
\item Environ 10GB par mois.
\item Environ 1GB par mois.
\item Environ 20GB par mois.
\end{enumerate}

\bigskip


\subsection*{Question 139}
De quoi traitait la guerre des blocs en 2017 ?

\begin{enumerate}[label=\Alph*., leftmargin=*]
\item Un conflit sur la question de savoir s'il fallait réduire la taille des blocs pour diminuer la capacité de transaction.
\item Un conflit sur la question de savoir s'il fallait diminuer la récompense de bloc pour décourager les mineurs.
\item Un conflit sur la question de savoir s'il fallait augmenter la taille des blocs pour étendre la capacité de transaction.
\item Un conflit sur la question de savoir s'il fallait augmenter la récompense de bloc pour inciter les mineurs.
\end{enumerate}

\bigskip


\subsection*{Question 140}
Pourquoi est-il important qu'un nœud Bitcoin reste abordable et accessible à tous les utilisateurs possibles ?

\begin{enumerate}[label=\Alph*., leftmargin=*]
\item Parce que cela facilite la décentralisation du réseau.
\item Parce que cela augmente la capacité de transaction du réseau.
\item Parce que cela diminue la capacité de transaction du réseau.
\item Parce que cela augmente la récompense de bloc pour les mineurs.
\end{enumerate}

\bigskip


\subsection*{Question 141}
Quelles seraient les implications pour les utilisateurs de Bitcoin si les blocs étaient 100 fois plus lourds ?

\begin{enumerate}[label=\Alph*., leftmargin=*]
\item La récompense de bloc pour les mineurs augmenterait.
\item La capacité de transaction du réseau diminuerait.
\item Faire fonctionner un nœud Bitcoin ne serait pas accessible à la personne moyenne, compromettant ainsi la décentralisation du protocole et l'immuabilité des transactions ainsi que des règles de consensus.
\item Faire fonctionner un nœud Bitcoin serait plus accessible à la personne moyenne, améliorant la décentralisation du protocole et l'immuabilité des transactions ainsi que des règles de consensus.
\end{enumerate}

\bigskip


\subsection*{Question 142}
Quel a été le résultat de la guerre des blocs en 2017 ?

\begin{enumerate}[label=\Alph*., leftmargin=*]
\item Les utilisateurs et les nœuds ont rejeté la proposition de modification visant à augmenter la taille des blocs, et ont activé une mise à jour appelée SegWit.
\item Les utilisateurs et les nœuds ont accepté la proposition de modification visant à augmenter la taille des blocs, et n'ont activé aucune mise à jour.
\item Les utilisateurs et les nœuds ont rejeté la proposition de modification visant à augmenter la taille des blocs, et n'ont activé aucune mise à jour.
\item Les utilisateurs et les nœuds ont accepté la proposition de modification visant à augmenter la taille des blocs, et ont activé une mise à jour appelée SegWit.
\end{enumerate}

\bigskip


\subsection*{Question 143}
Qu'est-ce que le Lightning Network ?

\begin{enumerate}[label=\Alph*., leftmargin=*]
\item Un réseau de mineurs Bitcoin.
\item Un réseau de paiement Bitcoin instantané basé sur la blockchain Bitcoin.
\item Un réseau blockchain séparé de la blockchain Bitcoin.
\item Un réseau d'échanges Bitcoin.
\end{enumerate}

\bigskip


\subsection*{Question 144}
Quel est le rôle des mineurs dans le réseau Bitcoin ?

\begin{enumerate}[label=\Alph*., leftmargin=*]
\item Les mineurs sécurisent le réseau et ajoutent des transactions aux blocs.
\item Les mineurs vérifient l'identité des utilisateurs de Bitcoin.
\item Les mineurs créent de nouveaux bitcoins.
\item Les mineurs fixent le prix du Bitcoin.
\end{enumerate}

\bigskip


\subsection*{Question 145}
Qu'est-ce que la Preuve de Travail (Proof of Work - POW) ?

\begin{enumerate}[label=\Alph*., leftmargin=*]
\item C'est l'algorithme utilisé pour miner des bitcoins.
\item C'est le consensus de sécurité du protocole Bitcoin.
\item C'est le processus de création de nouveaux bitcoins.
\item C'est la méthode de vérification des transactions Bitcoin.
\end{enumerate}

\bigskip


\subsection*{Question 146}
Qu'est-ce que le HashRate ?

\begin{enumerate}[label=\Alph*., leftmargin=*]
\item Le nombre total de mineurs dans le réseau Bitcoin.
\item La vitesse à laquelle un dispositif de minage peut résoudre la Preuve de Travail.
\item Le nombre de calculs de hachage qu'un dispositif de minage peut effectuer par seconde pour résoudre la Preuve de Travail.
\item La quantité totale de bitcoins minés en effectuant un certain nombre de tentatives par seconde pour résoudre la Preuve de Travail.
\end{enumerate}

\bigskip


\subsection*{Question 147}
Quel est le but de l'ajustement de difficulté dans le minage de Bitcoin ?

\begin{enumerate}[label=\Alph*., leftmargin=*]
\item Il augmente la sécurité du réseau Bitcoin.
\item Il augmente la récompense de minage.
\item Il diminue le nombre de bitcoins en circulation.
\item Il rééquilibre le jeu mondial de minage en fonction du nombre de participants.
\end{enumerate}

\bigskip


\subsection*{Question 148}
Quel est le rôle des ASICs dans le minage de Bitcoin ?

\begin{enumerate}[label=\Alph*., leftmargin=*]
\item Les ASICs sont des machines conçues uniquement pour appliquer l'algorithme SHA256 dans le processus de minage.
\item Les ASICs sont utilisés pour vérifier les transactions Bitcoin en appliquant l'algorithme SHA256 dans le processus de vérification.
\item Les ASICs sont utilisés pour créer arbitrairement de nouveaux bitcoins.
\item Les ASICs sont utilisés pour sécuriser le réseau Bitcoin.
\end{enumerate}

\bigskip


\subsection*{Question 149}
Qu'est-ce que la transaction coinbase ?

\begin{enumerate}[label=\Alph*., leftmargin=*]
\item La deuxième transaction du bloc qui inclut les frais les plus élevés.
\item Une transaction effectuée sur l'échange Coinbase.
\item Une transaction effectuée par le mineur après avoir reçu la récompense.
\item La première transaction du bloc qui inclut la récompense du mineur.
\end{enumerate}

\bigskip


\subsection*{Question 150}
Quel est le but des pools de minage ?

\begin{enumerate}[label=\Alph*., leftmargin=*]
\item Les pools de minage sont utilisés pour augmenter la récompense du minage.
\item Les pools de minage sont utilisés pour augmenter le HashRate du réseau Bitcoin dans son ensemble.
\item Les pools de minage permettent aux mineurs de regrouper leurs ressources informatiques et de recevoir des récompenses plus petites mais plus régulières.
\item Les pools de minage sont utilisés pour diminuer la difficulté du minage.
\end{enumerate}

\bigskip


\subsection*{Question 151}
Quel est le problème des généraux byzantins ?

\begin{enumerate}[label=\Alph*., leftmargin=*]
\item C'est un problème d'empêchement du double paiement dans une monnaie numérique.
\item C'est un problème de coordination des informations dans un système où les différents acteurs ne peuvent pas être considérés comme fiables.
\item C'est un problème de vérification des transactions dans un réseau décentralisé.
\item C'est un problème de maintien du consensus dans un réseau centralisé.
\end{enumerate}

\bigskip


\subsection*{Question 152}
Quel est le but du Proof of Work dans le réseau Bitcoin ?

\begin{enumerate}[label=\Alph*., leftmargin=*]
\item Le Proof of Work est utilisé pour vérifier les transactions Bitcoin en faisant confiance à un tiers.
\item Le Proof of Work est utilisé pour créer de nouveaux bitcoins en faisant confiance à un tiers.
\item Le Proof of Work est utilisé pour changer le protocole Bitcoin sans faire confiance à un tiers.
\item Le Proof of Work assure la véracité des informations sans faire confiance à un tiers.
\end{enumerate}

\bigskip


\subsection*{Question 153}
Quelle est l'importance de la dépense énergétique dans le minage de Bitcoin ?

\begin{enumerate}[label=\Alph*., leftmargin=*]
\item Le coût énergétique n'est pas nécessaire pour créer de nouveaux bitcoins.
\item La dépense énergétique est nécessaire pour maintenir le réseau Bitcoin, et la vérification de l'information a des coûts variables.
\item Le coût énergétique est nécessaire pour alimenter les machines ASIC utilisées dans le minage, et la vérification de l'information a des coûts élevés.
\item La dépense énergétique est nécessaire pour produire l'information, mais la vérification de l'information a un coût négligeable. Cette asymétrie garantit la sécurité du réseau.
\end{enumerate}

\bigskip


\subsection*{Question 154}
Que se passe-t-il lors d'une attaque à 51% sur le réseau Bitcoin ?

\begin{enumerate}[label=\Alph*., leftmargin=*]
\item Un agent prend le contrôle de plus de la moitié des pools de minage dans le réseau Bitcoin.
\item Un agent prend le contrôle de plus de la moitié des nœuds dans le réseau Bitcoin.
\item Un agent prend le contrôle de plus de la moitié des bitcoins en circulation.
\item Un agent possède plus de la moitié du hashrate et peut tenter de modifier la blockchain.
\end{enumerate}

\bigskip


\subsection*{Question 155}
Quelle est la principale préoccupation concernant le coût environnemental du minage de Bitcoin ?

\begin{enumerate}[label=\Alph*., leftmargin=*]
\item Le coût de l'électricité.
\item La consommation d'énergie.
\item Le recyclage des ASICs.
\item La production des ASICs.
\end{enumerate}

\bigskip


\subsection*{Question 156}
Quel est le rôle des mineurs dans le protocole Bitcoin ?

\begin{enumerate}[label=\Alph*., leftmargin=*]
\item Ils gèrent les transactions entre utilisateurs.
\item Ils contrôlent l'offre de Bitcoin.
\item Ils sécurisent le réseau actuel et historique.
\item Ils régulent le prix du Bitcoin.
\end{enumerate}

\bigskip


\subsection*{Question 157}
Quel est le temps moyen entre la création de deux blocs dans le protocole Bitcoin ?

\begin{enumerate}[label=\Alph*., leftmargin=*]
\item 15 minutes
\item 20 minutes
\item 5 minutes
\item 10 minutes
\end{enumerate}

\bigskip


\subsection*{Question 158}
Quel est le principal avantage de Bitcoin pour les individus vivant sous une oppression financière ou des régimes dictatoriaux ?

\begin{enumerate}[label=\Alph*., leftmargin=*]
\item Obtenir des rendements élevés sur l'investissement.
\item Avoir la possibilité d'échapper à la censure et aux restrictions bancaires.
\item Utiliser une monnaie stable qui a une valeur fixe.
\item Permettre des transactions anonymes.
\end{enumerate}

\bigskip


\subsection*{Question 159}
Comment le protocole Bitcoin maintient-il un temps moyen constant entre la création de deux blocs ?

\begin{enumerate}[label=\Alph*., leftmargin=*]
\item La récompense pour avoir trouvé un hash valide est augmentée.
\item La difficulté de trouver un hash valide s'ajuste tous les 2016 blocs.
\item La taille des blocs est ajustée à 10 mégaoctets.
\item Le nombre de mineurs est régulé à un nombre fixe.
\end{enumerate}

\bigskip


\subsection*{Question 160}
Que peuvent faire les mineurs pour les centrales électriques qui ne sont pas encore connectées au réseau ?

\begin{enumerate}[label=\Alph*., leftmargin=*]
\item Ils peuvent fournir de l'électricité aux centrales électriques.
\item Ils peuvent connecter les centrales électriques au réseau.
\item Ils peuvent agir comme acheteur de dernier recours.
\item Ils peuvent réguler le prix de l'électricité.
\end{enumerate}

\bigskip


\subsection*{Question 161}
Quel est le principal problème du système financier actuel en relation avec l'environnement ?

\begin{enumerate}[label=\Alph*., leftmargin=*]
\item Il favorise l'utilisation des combustibles fossiles.
\item Il encourage la surconsommation et l'endettement.
\item Il n'investit pas dans l'énergie renouvelable.
\item Il ne régule pas la production de biens.
\end{enumerate}

\bigskip


\subsection*{Question 162}
Comment le Bitcoin peut-il potentiellement aider à la transition écologique ?

\begin{enumerate}[label=\Alph*., leftmargin=*]
\item Il favorise l'utilisation d'énergies vertes.
\item Il réduit la production de biens.
\item Il régule la consommation de ressources.
\item Il limite l'émission uniquement de gaz à effet de serre.
\end{enumerate}

\bigskip


\subsection*{Question 163}
Comment le protocole Bitcoin récompense-t-il le mineur qui trouve un hash valide pour le bloc suivant ?

\begin{enumerate}[label=\Alph*., leftmargin=*]
\item Avec un pourcentage du total de l'offre de Bitcoin et les frais de toutes les transactions incluses dans les blocs précédents.
\item Avec les frais de toutes les transactions incluses dans le bloc valide.
\item Avec une quantité fixe de bitcoins non définie par les règles de consensus.
\item Avec une récompense définie par les règles de consensus, et avec les frais de toutes les transactions incluses dans le bloc valide.
\end{enumerate}

\bigskip


\subsection*{Question 164}
Quel est l'avantage principal des machines spécialisées SHA256, ou ASICs, dans le minage de Bitcoin ?

\begin{enumerate}[label=\Alph*., leftmargin=*]
\item Elles augmentent le taux de hachage par joule.
\item Elles réduisent la consommation d'énergie.
\item Elles augmentent la sécurité du réseau.
\item Elles augmentent la vitesse des transactions.
\end{enumerate}

\bigskip


\subsection*{Question 165}
Comment le protocole Bitcoin assure-t-il sa résistance à la censure et sa résilience ?

\begin{enumerate}[label=\Alph*., leftmargin=*]
\item Chaque composant du protocole est géographiquement distribué à travers le monde.
\item Il utilise un algorithme complexe pour valider les transactions.
\item Il utilise un réseau local de nœuds distribués en Europe.
\item Il utilise un haut niveau de cryptage.
\end{enumerate}

\bigskip


\subsection*{Question 166}
Quelle est la principale différence entre l'économie keynésienne et l'économie autrichienne concernant le système monétaire ?

\begin{enumerate}[label=\Alph*., leftmargin=*]
\item L'économie keynésienne promeut la surconsommation des monnaies.
\item L'économie keynésienne reconnaît la rareté des situations et des ressources.
\item L'économie keynésienne encourage l'inflation.
\item L'économie keynésienne ne prend pas en compte les aspects temporels et dynamiques des situations et des ressources.
\end{enumerate}

\bigskip


\subsection*{Question 167}
Quelle est l'une des raisons pour lesquelles le Bitcoin est attrayant pour certains groupes ?

\begin{enumerate}[label=\Alph*., leftmargin=*]
\item Il est contrôlé par une autorité centrale qui distribue librement du capital aux utilisateurs.
\item Il est facile à contrefaire, et peut créer des solutions pour recevoir plus d'argent d'un simple clic.
\item Il n'est pas pseudonyme, donc chaque utilisateur est bien connu et digne de confiance.
\item Sa capacité à offrir des solutions telles que l'argent électronique sans confiance, la résistance à la censure, ou une politique monétaire transparente et immuable.
\end{enumerate}

\bigskip


\subsection*{Question 168}
Quels sont les coûts typiques pour un mineur ?

\begin{enumerate}[label=\Alph*., leftmargin=*]
\item Consommation d'électricité et frais de transaction.
\item Achat de matériel uniquement.
\item Consommation d'électricité et achat de matériel.
\item Frais de transaction.
\end{enumerate}

\bigskip


\subsection*{Question 169}
Quelle est une manière d'atténuer la volatilité du Bitcoin ?

\begin{enumerate}[label=\Alph*., leftmargin=*]
\item Investir tout votre argent dans le Bitcoin.
\item Couvertures financières (stablecoins).
\item Ignorer les tendances du marché.
\item Acheter plus de bitcoin lorsque le prix baisse.
\end{enumerate}

\bigskip


\subsection*{Question 170}
Quelle est une caractéristique unique du protocole Bitcoin ?

\begin{enumerate}[label=\Alph*., leftmargin=*]
\item Il est contrôlé par une autorité centrale.
\item Il peut être facilement piraté.
\item Il fonctionne uniquement pendant les heures de bureau.
\item Il fonctionne à l'échelle mondiale, 24 heures sur 24, 7 jours sur 7.
\end{enumerate}

\bigskip


\subsection*{Question 171}
Quel a été un événement marquant dans le monde des cryptomonnaies en 2017 ?

\begin{enumerate}[label=\Alph*., leftmargin=*]
\item La disparition de toutes les cryptomonnaies.
\item Le lancement de milliers d'Offres Initiales de Pièces (ICO).
\item La création du Bitcoin.
\item L'interdiction du Bitcoin.
\end{enumerate}

\bigskip


\subsection*{Question 172}
Quelle est la tendance générale du Bitcoin, compte tenu de son offre limitée et du processus de halving ?

\begin{enumerate}[label=\Alph*., leftmargin=*]
\item Une tendance générale à la baisse.
\item Une tendance générale à la hausse.
\item Aucune tendance particulière.
\item Une tendance qui fluctue de manière extrême sans motif discernable.
\end{enumerate}

\bigskip


\subsection*{Question 173}
Quelle est l'une des raisons pour lesquelles le Bitcoin connaît tant de bulles spéculatives ?

\begin{enumerate}[label=\Alph*., leftmargin=*]
\item Parce que le Bitcoin est une bulle qui va disparaître.
\item Parce que le Bitcoin est une monnaie qui est utilisée partout dans le monde.
\item Parce que le Bitcoin n'est pas soumis aux cycles économiques.
\item Parce que le Bitcoin n'est réellement utilisé nulle part.
\end{enumerate}

\bigskip


\subsection*{Question 174}
Quel est un facteur qui rend difficile pour les autorités financières de réguler le Bitcoin ?

\begin{enumerate}[label=\Alph*., leftmargin=*]
\item Il est facile à contrefaire.
\item Il n'est utilisé nulle part dans le monde.
\item Il est contrôlé par une autorité centrale.
\item Il fonctionne à l'échelle mondiale, 24 heures sur 24, 7 jours sur 7.
\end{enumerate}

\bigskip


\subsection*{Question 175}
Quelle est une manière dont le Bitcoin s'est davantage intégré au marché traditionnel ?

\begin{enumerate}[label=\Alph*., leftmargin=*]
\item La création du Bitcoin.
\item L'arrivée des ETF Bitcoin.
\item La disparition de toutes les cryptomonnaies.
\item L'interdiction du Bitcoin.
\end{enumerate}

\bigskip


\subsection*{Question 176}
Parmi les caractéristiques suivantes, laquelle a contribué à l'expansion de l'utilisation de Bitcoin dans les marchés noirs du web, comme Silk Road ?

\begin{enumerate}[label=\Alph*., leftmargin=*]
\item Sa nature incontrôlable et pseudonyme.
\item Son absence de pseudonymat.
\item Sa traçabilité et sa nature contrôlable.
\item Son contrôle centralisé.
\end{enumerate}

\bigskip


\subsection*{Question 177}
Quelle est une raison pour laquelle le prix du Bitcoin peut chuter de 10 %, 20 %, ou même 50 % en seulement quelques jours ?

\begin{enumerate}[label=\Alph*., leftmargin=*]
\item La météo.
\item Le nombre de personnes utilisant internet.
\item Les phases haussières et baissières du marché.
\item Le prix de l'or.
\end{enumerate}

\bigskip


\subsection*{Question 178}
Quelle est l'une des raisons pour lesquelles le Bitcoin n'est pas une bulle qui va disparaître ?

\begin{enumerate}[label=\Alph*., leftmargin=*]
\item Parce que c'est une monnaie qui est réellement utilisée partout dans le monde.
\item Parce qu'elle est facile à contrefaire.
\item Parce qu'elle est contrôlée par une autorité centrale.
\item Parce qu'elle n'est réellement utilisée nulle part.
\end{enumerate}

\bigskip


\subsection*{Question 179}
En comparaison avec les monnaies fiduciaires, comment est perçu le Bitcoin ?

\begin{enumerate}[label=\Alph*., leftmargin=*]
\item Comme une économie tertiaire.
\item Comme une économie parallèle.
\item Comme une économie de remplacement.
\item Comme une économie industrielle.
\end{enumerate}

\bigskip


\subsection*{Question 180}
Quelle est une manière d'obtenir des bitcoins sans les acheter ?

\begin{enumerate}[label=\Alph*., leftmargin=*]
\item Les trouver dans la rue.
\item Les gagner à la loterie.
\item Les accepter comme moyen de paiement pour les produits ou services que vous proposez.
\item Les emprunter à un ami et ne jamais les rendre.
\end{enumerate}

\bigskip


\subsection*{Question 181}
Quel est un avantage d'accepter le Bitcoin en tant que commerçant ?

\begin{enumerate}[label=\Alph*., leftmargin=*]
\item Résistance à la censure.
\item Restriction financière.
\item Augmentation des frais de transaction.
\item Protection contre la déflation.
\end{enumerate}

\bigskip


\subsection*{Question 182}
Qu'est-ce que BTCMap ?

\begin{enumerate}[label=\Alph*., leftmargin=*]
\item Une carte de toutes les mines de Bitcoin dans le monde.
\item Une plateforme qui répertorie tous les commerçants qui acceptent le Bitcoin.
\item Une carte de tous les échanges de Bitcoin dans le monde.
\item Une carte de tous les utilisateurs de Bitcoin dans le monde.
\end{enumerate}

\bigskip


\subsection*{Question 183}
Comment Bitcoin résout-il le problème des généraux byzantins ?

\begin{enumerate}[label=\Alph*., leftmargin=*]
\item Par l'utilisation d'une blockchain basée sur la Preuve de Travail avec une difficulté variable.
\item Par l'utilisation de signatures numériques.
\item Par l'utilisation d'un registre public.
\item Par l'utilisation d'un algorithme de Preuve d'Enjeu.
\end{enumerate}

\bigskip


\subsection*{Question 184}
Qu'est-ce que l'AMF ?

\begin{enumerate}[label=\Alph*., leftmargin=*]
\item Autorité des Marchés Financiers.
\item Fonds Monétaire Américain.
\item Association des Installations Minières.
\item Flux Monétaire Automatisé.
\end{enumerate}

\bigskip


\subsection*{Question 185}
Quel est un facteur à considérer lors du choix d'une solution pour accepter le Bitcoin pour votre entreprise ?

\begin{enumerate}[label=\Alph*., leftmargin=*]
\item L'emplacement de votre siège social.
\item Le volume de transactions attendu.
\item Le nombre d'employés.
\item La couleur de votre logo.
\end{enumerate}

\bigskip


\subsection*{Question 186}
Qu'est-ce que OpenNode ?

\begin{enumerate}[label=\Alph*., leftmargin=*]
\item Une solution en ligne simple pour que les entreprises acceptent le Bitcoin.
\item Un type de nœud Bitcoin qui est utilisé simultanément par de nombreux utilisateurs.
\item Une plateforme d'échange de Bitcoin.
\item Un logiciel de minage de Bitcoin qui est open-source.
\end{enumerate}

\bigskip


\subsection*{Question 187}
Quelle est l'une des initiatives qui contribuent à rendre l'économie Bitcoin plus accessible et pratique pour tous ?

\begin{enumerate}[label=\Alph*., leftmargin=*]
\item Loterie Bitcoin.
\item Minage Bitcoin.
\item Échange Bitcoin.
\item BTCMap.
\end{enumerate}

\bigskip


\subsection*{Question 188}
Quel est un outil permettant aux grandes organisations ou aux passionnés de Bitcoin d'accepter le Bitcoin ?

\begin{enumerate}[label=\Alph*., leftmargin=*]
\item Échange Bitcoin
\item Portefeuille Bitcoin
\item BTCpay Server
\item Minage BTCpay
\end{enumerate}

\bigskip


\subsection*{Question 189}
Parmi les actions suivantes, lesquelles peuvent aider ceux qui souhaitent utiliser Bitcoin dans la vie quotidienne ?

\begin{enumerate}[label=\Alph*., leftmargin=*]
\item Diminuer le nombre de bitcoins.
\item Limiter l'utilisation du Bitcoin.
\item Avoir une liste de tous les commerçants qui acceptent Bitcoin.
\item Augmenter le prix du Bitcoin.
\end{enumerate}

\bigskip


\subsection*{Question 190}
Qu'est-ce que Swiss Bitcoin Pay ?

\begin{enumerate}[label=\Alph*., leftmargin=*]
\item Une solution pour les commerçants amateurs d'accepter le Bitcoin.
\item Une plateforme d'échange de Bitcoin en Suisse.
\item Une banque suisse qui accepte le Bitcoin.
\item Une entreprise minière qui opère en Suisse.
\end{enumerate}

\bigskip


\subsection*{Question 191}
Quel est l'un des risques associés à l'achat de Bitcoin ?

\begin{enumerate}[label=\Alph*., leftmargin=*]
\item Son prix peut chuter à 0.
\item Il peut être dupliqué.
\item Il peut être perdu physiquement.
\item Il peut être facilement volé.
\end{enumerate}

\bigskip


\subsection*{Question 192}
Quelle est l'une des choses que vous devriez avoir avant d'acheter du Bitcoin ?

\begin{enumerate}[label=\Alph*., leftmargin=*]
\item Une carte de crédit.
\item Un portefeuille sécurisé.
\item Un compte en banque.
\item Un coffre-fort physique.
\end{enumerate}

\bigskip


\subsection*{Question 193}
Qu'est-ce que le Dollar Cost Average ?

\begin{enumerate}[label=\Alph*., leftmargin=*]
\item Acheter du bitcoin uniquement lorsque le prix baisse.
\item Acheter une grande quantité de bitcoin en une seule fois.
\item Acheter de petites quantités de bitcoin à intervalles réguliers.
\item Acheter du bitcoin uniquement lorsque le prix monte.
\end{enumerate}

\bigskip


\subsection*{Question 194}
Qu'est-ce qu'un achat spontané ?

\begin{enumerate}[label=\Alph*., leftmargin=*]
\item Acheter du bitcoin à intervalles réguliers.
\item Acheter du bitcoin uniquement lorsque le prix monte.
\item Acheter du bitcoin uniquement lorsque le prix baisse.
\item Acheter immédiatement du bitcoin pour obtenir une exposition.
\end{enumerate}

\bigskip


\subsection*{Question 195}
Quel est le but des réglementations KYC ?

\begin{enumerate}[label=\Alph*., leftmargin=*]
\item Contrôler l'offre de Bitcoin.
\item Combattre le financement du terrorisme, l'évasion fiscale et le blanchiment d'argent.
\item Empêcher les gens d'acheter du Bitcoin.
\item Augmenter le prix du Bitcoin et faire payer plus d'impôts aux utilisateurs.
\end{enumerate}

\bigskip


\subsection*{Question 196}
Quelle est l'une des manières d'acquérir des bitcoins sans KYC ?

\begin{enumerate}[label=\Alph*., leftmargin=*]
\item Plateformes de courtage
\item Transferts bancaires
\item Plateformes DCA
\item Distributeurs automatiques de Bitcoin
\end{enumerate}

\bigskip


\subsection*{Question 197}
Que devez-vous faire après avoir acheté des bitcoins sur une plateforme d'échange ?

\begin{enumerate}[label=\Alph*., leftmargin=*]
\item Retirer les bitcoins et les envoyer à votre portefeuille.
\item Transférer les bitcoins vers une autre plateforme d'échange.
\item Vendre immédiatement les bitcoins.
\item Laisser les bitcoins sur la plateforme.
\end{enumerate}

\bigskip


\subsection*{Question 198}
Quel est l'un des éléments à prendre en compte lors du choix d'une plateforme DCA ?

\begin{enumerate}[label=\Alph*., leftmargin=*]
\item La couleur du logo de la plateforme.
\item L'âge de la plateforme.
\item Le nombre d'utilisateurs sur la plateforme.
\item La disponibilité dans votre pays.
\end{enumerate}

\bigskip


\subsection*{Question 199}
Quelle est l'une des raisons pour lesquelles vous devriez consolider vos UTXOs dans vos portefeuilles de temps en temps ?

\begin{enumerate}[label=\Alph*., leftmargin=*]
\item Pour augmenter la valeur de vos bitcoins.
\item Pour empêcher vos bitcoins d'être volés.
\item Pour éviter des frais inutiles lors des transactions.
\item Pour garantir que vos bitcoins sont toujours disponibles.
\end{enumerate}

\bigskip


\subsection*{Question 200}
Que sont le FOMO et le FUD ?

\begin{enumerate}[label=\Alph*., leftmargin=*]
\item Des réglementations liées au Bitcoin.
\item Des types de portefeuilles Bitcoin.
\item Des stratégies pour acheter du Bitcoin.
\item Des émotions qui peuvent conduire à des prises de décision impulsives et potentiellement nocives.
\end{enumerate}

\bigskip


\subsection*{Question 201}
Quel est l'un des risques potentiels de laisser vos bitcoins sur une plateforme d'échange ?

\begin{enumerate}[label=\Alph*., leftmargin=*]
\item Le risque que la plateforme contrôle l'offre de Bitcoin.
\item Le risque que le Bitcoin soit dépassé en prix par une autre cryptomonnaie.
\item Le risque que la plateforme augmente le prix du Bitcoin.
\item Les risques de piratage et de blocage des fonds.
\end{enumerate}

\bigskip


\subsection*{Question 202}
Quelle est l'une des conséquences potentielles de ne pas suivre les réglementations de votre pays lors de l'achat de Bitcoin ?

\begin{enumerate}[label=\Alph*., leftmargin=*]
\item Vous pourriez vous mettre en danger.
\item Vous pourriez augmenter le prix du Bitcoin.
\item Vous pourriez contrôler l'offre de Bitcoin.
\item Vous pourriez empêcher les autres d'acheter du Bitcoin.
\end{enumerate}

\bigskip


\subsection*{Question 203}
À quoi ressemble la courbe d'adoption du Bitcoin ?

\begin{enumerate}[label=\Alph*., leftmargin=*]
\item Une courbe en W.
\item Une courbe en U.
\item Une courbe en V.
\item Une courbe en S.
\end{enumerate}

\bigskip


\subsection*{Question 204}
Quel pays a déclaré le Bitcoin comme monnaie légale ?

\begin{enumerate}[label=\Alph*., leftmargin=*]
\item Inde
\item El Salvador
\item Chine
\item États-Unis
\end{enumerate}

\bigskip


\subsection*{Question 205}
Quelle action Bitcoin force-t-il les entreprises, les universités, les régulateurs et les individus à entreprendre ?

\begin{enumerate}[label=\Alph*., leftmargin=*]
\item À l'ignorer.
\item À le criminaliser.
\item À le prendre en compte.
\item À l'interdire.
\end{enumerate}

\bigskip


\subsection*{Question 206}
Comment ce cours décrit-il Bitcoin ?

\begin{enumerate}[label=\Alph*., leftmargin=*]
\item Un sujet physique.
\item Un sujet simple.
\item Un sujet mono-disciplinaire.
\item Un sujet multidisciplinaire.
\end{enumerate}

\bigskip


\subsection*{Question 207}
Quel est l'objectif principal de la révolution Bitcoin selon ce cours ?

\begin{enumerate}[label=\Alph*., leftmargin=*]
\item Créer un monde où l'argent physique est obsolète.
\item Créer un monde où la vie privée n'est pas un droit humain, et l'argent est manipulé par les banques.
\item Créer un monde où les gouvernements n'ont aucun pouvoir, et la société est chaotique.
\item Créer un monde où la souveraineté humaine est un droit, la vie privée est respectée par défaut, et l'argent n'est pas manipulé.
\end{enumerate}

\bigskip


\subsection*{Question 208}
Qu'a prédit Milton Friedman en 1999 ?

\begin{enumerate}[label=\Alph*., leftmargin=*]
\item La disparition de l'argent physique.
\item Le développement d'un e-cash fiable.
\item La fin du contrôle gouvernemental sur l'argent.
\item L'émergence de Bitcoin.
\end{enumerate}

\bigskip


\subsection*{Question 209}
Que signifie le fait que Bitcoin ait une tendance générale à la hausse ?

\begin{enumerate}[label=\Alph*., leftmargin=*]
\item Une tendance générale à la hausse pour Bitcoin signifie que, globalement, le prix du Bitcoin a été en augmentation.
\item Une tendance générale à la hausse pour Bitcoin signifie que le prix est garanti de monter indéfiniment sans aucune fluctuation.
\item Une tendance générale à la hausse pour Bitcoin signifie qu'il est seulement populaire parmi un petit groupe d'investisseurs et n'a aucun impact sur le marché global.
\item Une tendance générale à la hausse pour Bitcoin signifie que Bitcoin devient moins précieux et perd son attrait sur le marché.
\end{enumerate}

\bigskip


\subsection*{Question 210}
Quel est le statut du Bitcoin au Salvador ?

\begin{enumerate}[label=\Alph*., leftmargin=*]
\item Inconnu
\item Criminalisé
\item Non réglementé
\item Monnaie légale
\end{enumerate}

\bigskip


\subsection*{Question 211}
Dans quels domaines le Bitcoin soulève-t-il des questions ?

\begin{enumerate}[label=\Alph*., leftmargin=*]
\item Cryptographie, théorie des jeux, économie et politique monétaire, informatique, philosophie, énergie, lois et réglementation.
\item Uniquement la philosophie et l'énergie.
\item Uniquement la cryptographie et la théorie des jeux.
\item Uniquement l'informatique et l'économie.
\end{enumerate}

\bigskip


\subsection*{Question 212}
Quelle est la nature de l'adoption de Bitcoin décrite dans ce cours ?

\begin{enumerate}[label=\Alph*., leftmargin=*]
\item Virale
\item Rapide et incontrôlable
\item Lente et régulière
\item Sporadique et imprévisible
\end{enumerate}

\bigskip


\subsection*{Question 213}
Que suggère ce cours sur l'avenir du Bitcoin ?

\begin{enumerate}[label=\Alph*., leftmargin=*]
\item Il perdra sa valeur.
\item Il sera remplacé par une meilleure technologie.
\item Il ne va pas disparaître.
\item Il sera interdit dans le monde entier.
\end{enumerate}

\bigskip


\subsection*{Question 214}
Que suggère ce cours sur le rythme d'apprentissage concernant Bitcoin ?

\begin{enumerate}[label=\Alph*., leftmargin=*]
\item Il est facile d'apprendre tout d'un coup.
\item Il est impossible d'apprendre tout sur Bitcoin.
\item Il est compliqué d'assimiler tout d'un coup, donc prenez votre temps.
\item Il n'est pas nécessaire d'apprendre tout sur Bitcoin.
\end{enumerate}

\bigskip


\subsection*{Question 215}
Qu'est-ce que le Lightning Network ?

\begin{enumerate}[label=\Alph*., leftmargin=*]
\item Une méthode pour miner du Bitcoin.
\item Une nouvelle technologie blockchain qui remplacera Bitcoin à l'avenir.
\item Un réseau de paiement qui utilise le protocole Bitcoin pour permettre des transactions ultra-rapides.
\item Un nouveau type de cryptomonnaie.
\end{enumerate}

\bigskip


\subsection*{Question 216}
Quel est le problème de scalabilité de Bitcoin ?

\begin{enumerate}[label=\Alph*., leftmargin=*]
\item Le manque de régulation pour Bitcoin qui rendra les transactions illégales.
\item La difficulté de miner de nouveaux bitcoins avec une valeur augmentant au fil du temps.
\item La volatilité du prix du Bitcoin qui augmentera après chaque halving.
\item Le défi de mettre en place un système monétaire capable de fournir un nombre toujours croissant de transactions par seconde à mesure qu'il est adopté.
\end{enumerate}

\bigskip


\subsection*{Question 217}
Qu'est-ce que le trilemme de la blockchain ?

\begin{enumerate}[label=\Alph*., leftmargin=*]
\item Un protocole basé sur une blockchain ne peut satisfaire que deux des critères parmi la décentralisation, la sécurité et la scalabilité.
\item Un protocole basé sur une blockchain peut satisfaire tous les critères de décentralisation, de sécurité et de scalabilité.
\item Un protocole basé sur une blockchain ne peut satisfaire aucun des critères de décentralisation, de sécurité et de scalabilité.
\item Un protocole basé sur une blockchain ne peut satisfaire qu'un seul des critères parmi la décentralisation, la sécurité et la scalabilité.
\end{enumerate}

\bigskip


\subsection*{Question 218}
Que permet le Lightning Network ?

\begin{enumerate}[label=\Alph*., leftmargin=*]
\item Des transactions Bitcoin retardées et à coût élevé.
\item Des transactions Bitcoin retardées et à faible coût.
\item Des transactions Bitcoin instantanées et à coût élevé.
\item Des transactions Bitcoin instantanées et à faible coût.
\end{enumerate}

\bigskip


\subsection*{Question 219}
Quand le Lightning Network a-t-il été validé et implémenté pour la première fois ?

\begin{enumerate}[label=\Alph*., leftmargin=*]
\item 2015
\item 2018
\item 2016
\item 2017
\end{enumerate}

\bigskip


\subsection*{Question 220}
Quel type de réseau le Lightning Network crée-t-il ?

\begin{enumerate}[label=\Alph*., leftmargin=*]
\item Un réseau d'échanges de cryptomonnaies.
\item Un réseau de canaux de paiement.
\item Un réseau de nœuds de minage.
\item Un réseau de nœuds de blockchain.
\end{enumerate}

\bigskip


\subsection*{Question 221}
Qu'arriverait-il aux services traditionnels de transfert d'argent comme Western Union, aux banques centrales, à Visa et à Mastercard s'ils n'adoptaient pas le Lightning Network ?

\begin{enumerate}[label=\Alph*., leftmargin=*]
\item Ils pourraient devenir plus décentralisés.
\item Ils pourraient devenir plus populaires.
\item Ils pourraient devenir plus sécurisés.
\item Ils pourraient disparaître.
\end{enumerate}

\bigskip


\subsection*{Question 222}
Les frais pour les transactions sur le Lightning Network sont extrêmement bas, souvent inférieurs à... ?

\begin{enumerate}[label=\Alph*., leftmargin=*]
\item 3 %
\item 0,5 %
\item 2 %
\item 1 %
\end{enumerate}

\bigskip


\subsection*{Question 223}
Comment les transactions sur le Lightning Network sont-elles sécurisées ?

\begin{enumerate}[label=\Alph*., leftmargin=*]
\item Par l'utilisation uniquement des deux, clés privées et publiques.
\item Par la cryptographie et indirectement par l'énergie consommée par les mineurs sur Bitcoin.
\item Par l'utilisation uniquement de clés privées.
\item Par l'utilisation uniquement de clés publiques.
\end{enumerate}

\bigskip


\subsection*{Question 224}
Qu'est-ce que le Lightning Network rend possible ?

\begin{enumerate}[label=\Alph*., leftmargin=*]
\item Établir des connexions de canal direct qui sont limitées uniquement aux grandes transactions.
\item Réaliser des transactions entre individus qui n'ont pas de connexion directe de canal.
\item Créer des connexions de canal qui nécessitent plusieurs confirmations pour toute transaction.
\item Créer des connexions de canal direct uniquement avec des échanges centralisés.
\end{enumerate}

\bigskip


\subsection*{Question 225}
Quel est l'intervalle de temps moyen entre la création de deux blocs dans le protocole Bitcoin ?

\begin{enumerate}[label=\Alph*., leftmargin=*]
\item 5 minutes
\item 20 minutes
\item 10 minutes
\item 15 minutes
\end{enumerate}

\bigskip


\subsection*{Question 226}
Quelle est la taille d'un bloc dans le protocole Bitcoin ?

\begin{enumerate}[label=\Alph*., leftmargin=*]
\item 2MB.
\item 4MB.
\item 1MB.
\item 3MB.
\end{enumerate}

\bigskip


\subsection*{Question 227}
Quel est le but principal du Lightning Network ?

\begin{enumerate}[label=\Alph*., leftmargin=*]
\item Créer une nouvelle cryptomonnaie.
\item Remplacer le protocole Bitcoin.
\item Augmenter la valeur du Bitcoin.
\item Faciliter les micro-transactions.
\end{enumerate}

\bigskip


\subsection*{Question 228}
Quelle est l'une des applications potentielles du Lightning Network dans la vie quotidienne ?

\begin{enumerate}[label=\Alph*., leftmargin=*]
\item Créer de nouvelles cryptomonnaies.
\item Augmenter la valeur du Bitcoin.
\item Miner de nouveaux Bitcoins.
\item Facturation instantanée dans le commerce.
\end{enumerate}

\bigskip


\subsection*{Question 229}
Quelle est l'unité de Bitcoin utilisée dans le Lightning Network ?

\begin{enumerate}[label=\Alph*., leftmargin=*]
\item Block
\item Ether
\item Litecoin
\item Sats
\end{enumerate}

\bigskip


\subsection*{Question 230}
Quel est le concept de skin in the game dans le contexte de l'industrie du jeu vidéo ?

\begin{enumerate}[label=\Alph*., leftmargin=*]
\item Vendre des skins de jeu pour de l'argent réel.
\item Personnaliser des personnages avec des skins uniques qui améliorent le gameplay.
\item Acheter des objets dans le jeu avec de l'argent réel.
\item Avoir un enjeu financier dans le résultat d'un jeu.
\end{enumerate}

\bigskip


\subsection*{Question 231}
Comment le Lightning Network peut-il potentiellement transformer les modèles d'affaires existants dans l'industrie du jeu vidéo ?

\begin{enumerate}[label=\Alph*., leftmargin=*]
\item Il permet aux joueurs d'acheter des objets en jeu avec du Bitcoin.
\item Il permet aux développeurs de jeux d'accepter le Bitcoin comme moyen de paiement pour leurs jeux.
\item Il permet aux joueurs de miser de très petites sommes d'argent lorsqu'ils jouent à des jeux.
\item Il permet aux joueurs de miner du Bitcoin en jouant.
\end{enumerate}

\bigskip


\subsection*{Question 232}
Quel est le concept de money streaming dans le contexte du Lightning Network ?

\begin{enumerate}[label=\Alph*., leftmargin=*]
\item Convertir des Bitcoins en espèces dans un emplacement physique à l'aide d'un distributeur automatique.
\item Effectuer des micropaiements continus, en temps réel.
\item Créer de nouveaux bitcoins par des opérations de minage qui se déroulent en temps réel.
\item Diffuser du contenu audio ou vidéo pour des bitcoins.
\end{enumerate}

\bigskip


\subsection*{Question 233}
Lequel des éléments suivants n'est pas une caractéristique clé des micro-transactions sur le Lightning Network ?

\begin{enumerate}[label=\Alph*., leftmargin=*]
\item Temps de paiement longs
\item Paiements instantanés
\item Frais faibles
\item Scalabilité
\end{enumerate}

\bigskip


\subsection*{Question 234}
Quel est l'un des avantages potentiels du Lightning Network pour les créateurs de contenu ?

\begin{enumerate}[label=\Alph*., leftmargin=*]
\item Ils pourraient facturer plus cher pour leur contenu.
\item Ils seraient incités à produire du contenu de bonne qualité pour retenir l'attention des utilisateurs.
\item Ils pourraient toucher un public plus large.
\item Ils pourraient produire du contenu plus rapidement pour retenir l'attention des utilisateurs.
\end{enumerate}

\bigskip


\subsection*{Question 235}
Qu'est-ce qu'une plateforme DCA ?

\begin{enumerate}[label=\Alph*., leftmargin=*]
\item Un outil logiciel qui suit les prix des actions en temps réel sans aucune fonctionnalité d'investissement.
\item Un service qui permet aux utilisateurs d'investir un montant fixe d'argent dans un actif à intervalles réguliers.
\item Un type de bourse de cryptomonnaies qui permet uniquement le trading de Bitcoin.
\item Un service financier qui fournit des prêts aux investisseurs pour le trading sur marge.
\end{enumerate}

\bigskip


\subsection*{Question 236}
Comment le réseau Lightning pourrait-il potentiellement être utilisé pour la location de biens ?

\begin{enumerate}[label=\Alph*., leftmargin=*]
\item Les utilisateurs pourraient être facturés en fonction de la valeur des biens.
\item Les utilisateurs pourraient être facturés en fonction de l'état des biens.
\item Les utilisateurs pourraient être facturés à la minute, ou même à la seconde, pour le temps d'utilisation des biens.
\item Les utilisateurs pourraient être facturés un tarif élevé pour la location des biens.
\end{enumerate}

\bigskip


\subsection*{Question 237}
Quelles sont certaines opportunités commerciales offertes par le Lightning Network ?

\begin{enumerate}[label=\Alph*., leftmargin=*]
\item Nouveaux modèles économiques où les consommateurs ne paient pas pour aucun contenu.
\item Nouveaux modèles économiques où les consommateurs utilisent des blockchains personnalisées.
\item Nouveaux modèles économiques où les consommateurs paient pour le contenu en fonction de ce qu'ils consomment.
\item Nouveaux modèles économiques où les consommateurs paient en utilisant différentes cryptomonnaies.
\end{enumerate}

\bigskip


\subsection*{Question 238}
Quelle est une manière de tester le Lightning Network par vous-même ?

\begin{enumerate}[label=\Alph*., leftmargin=*]
\item En minant du Bitcoin avec un ordinateur personnel.
\item En essayant l'application de podcast appelée Fountain.
\item En créant une nouvelle cryptomonnaie.
\item En achetant du Bitcoin auprès d'une banque.
\end{enumerate}

\bigskip


\subsection*{Question 239}
Quelle est une implication potentielle de l'avancement des technologies IA ?

\begin{enumerate}[label=\Alph*., leftmargin=*]
\item Les technologies IA entraîneront la disparition de plus de 80 % des emplois actuels.
\item Les technologies IA détruiront tous les ordinateurs.
\item Les technologies IA provoqueront un effondrement économique mondial.
\item Les technologies IA entraîneront une diminution de l'intelligence humaine.
\end{enumerate}

\bigskip


\subsection*{Question 240}
Quel est le lien entre Bitcoin et l'avenir de la technologie et de la société ?

\begin{enumerate}[label=\Alph*., leftmargin=*]
\item Bitcoin est une forme de technologie IA qui permet aux gens d'échanger des informations via un serveur central.
\item Bitcoin est un outil pour contrôler la population mondiale, permettant l'élimination de la pauvreté.
\item Bitcoin est une méthode pour que les gouvernements suivent les transactions financières de leurs citoyens.
\item Bitcoin est une révolution technologique qui permet l'échange de valeur sans aucun tiers de confiance.
\end{enumerate}

\bigskip


\subsection*{Question 241}
Quelle est l'une des principales questions soulevées concernant l'avenir de la finance ?

\begin{enumerate}[label=\Alph*., leftmargin=*]
\item Comment pouvons-nous garantir que tout le monde ait la même quantité d'argent ?
\item Qui devrait détenir, autoriser et tracer l'argent que nous utilisons ?
\item Comment pouvons-nous éliminer complètement l'argent ?
\item Comment pouvons-nous rendre l'argent obsolète ?
\end{enumerate}

\bigskip


\subsection*{Question 242}
Quelle est l'une des principales caractéristiques du Bitcoin ?

\begin{enumerate}[label=\Alph*., leftmargin=*]
\item Il ne peut être utilisé que par un groupe sélectionné de personnes.
\item Il nécessite une autorisation d'une autorité centrale pour être utilisé.
\item Il peut être échangé partout dans le monde sans aucune autorisation d'une quelconque entité.
\item Il ne peut être utilisé que dans certains pays.
\end{enumerate}

\bigskip


\subsection*{Question 243}
Quel est l'un des avantages potentiels de l'acceptation de nouvelles technologies telles que Bitcoin ?

\begin{enumerate}[label=\Alph*., leftmargin=*]
\item Cela pourrait conduire à l'élimination de tous les emplois.
\item Cela pourrait conduire à un effondrement économique mondial.
\item Cela pourrait générer d'énormes économies à l'échelle mondiale.
\item Cela pourrait entraîner une diminution de l'intelligence humaine.
\end{enumerate}

\bigskip


\subsection*{Question 244}
Quelle pourrait être une préoccupation majeure concernant le système bancaire ?

\begin{enumerate}[label=\Alph*., leftmargin=*]
\item Le système bancaire est trop facile à comprendre.
\item Les règles du jeu bancaire ne sont pas transparentes et compréhensibles pour tous.
\item Le système bancaire est trop transparent.
\item Le système bancaire est trop décentralisé.
\end{enumerate}

\bigskip


\subsection*{Question 245}
Quel est l'un des principaux problèmes liés à la censure ?

\begin{enumerate}[label=\Alph*., leftmargin=*]
\item Elle peut entraîner une diminution de l'innovation.
\item Elle peut conduire à davantage de soutien pour le droit de réunion.
\item Elle peut détruire la liberté d'expression et le droit de réunion.
\item Elle peut conduire à l'application de la liberté d'expression.
\end{enumerate}

\bigskip


\subsection*{Question 246}
Quel est l'un des principaux avantages du Bitcoin pour les personnes sans compte bancaire ?

\begin{enumerate}[label=\Alph*., leftmargin=*]
\item Le Bitcoin bénéficie uniquement aux personnes ayant un statut social élevé.
\item Le Bitcoin favorise l'inégalité dans les transactions, en récompensant les utilisateurs les plus riches.
\item Le Bitcoin bénéficie uniquement aux personnes ayant un compte bancaire et une position politique de gauche.
\item Le Bitcoin permet une égalité dans les transactions, indépendamment de votre statut social ou de votre position politique.
\end{enumerate}

\bigskip


\subsection*{Question 247}
Quelle est l'une des principales caractéristiques du protocole Bitcoin ?

\begin{enumerate}[label=\Alph*., leftmargin=*]
\item Il est neutre envers tous ses utilisateurs.
\item Il bénéficie uniquement à un groupe sélectionné d'utilisateurs.
\item Il favorise certains utilisateurs par rapport à d'autres.
\item Il est biaisé envers certains utilisateurs.
\end{enumerate}

\bigskip


\subsection*{Question 248}
Quel est un moyen pratique d'obtenir des bitcoins ?

\begin{enumerate}[label=\Alph*., leftmargin=*]
\item Les miner avec un ordinateur portable.
\item Les trouver dans le monde physique.
\item Les créer en utilisant un logiciel spécial.
\item Les acheter via des plateformes régulées ou non régulées.
\end{enumerate}

\bigskip


\subsection*{Question 249}
Quelle est l'une des principales raisons de la création de Bitcoin ?

\begin{enumerate}[label=\Alph*., leftmargin=*]
\item Proposer le contrôle de la population mondiale.
\item Proposer l'élimination de l'argent.
\item Proposer le changement du système financier.
\item Proposer la création d'une nouvelle forme de technologie IA.
\end{enumerate}

\bigskip


\subsection*{Question 250}
Quelle est l'une des principales conséquences du Bitcoin sur le système bancaire ?

\begin{enumerate}[label=\Alph*., leftmargin=*]
\item Il permet la corruption du système bancaire.
\item Il permet le contrôle du système bancaire.
\item Il permet l'émancipation du système bancaire.
\item Il permet l'élimination du système bancaire.
\end{enumerate}

\bigskip


\newpage
\section*{Answer Key}

\begin{enumerate}
\item \textbf{Question 1:} D
\\ \textit{La guerre de la taille des blocs a eu lieu en 2017. Ce conflit opposait les acteurs qui voulaient modifier Bitcoin en augmentant la taille des blocs pour augmenter la capacité des transactions, à ceux qui cherchaient à préserver l'indépendance et le pouvoir des utilisateurs en maintenant la taille à 1 MB.}

\item \textbf{Question 2:} D
\\ \textit{Le but de ce cours sur Bitcoin est de discuter des caractéristiques monétaires générales de Bitcoin, y compris comment acheter et vendre des bitcoins, les stocker de manière sécurisée dans des portefeuilles numériques, et les utiliser pour des transactions. Il vise également à explorer l'avenir de Bitcoin.}

\item \textbf{Question 3:} B
\\ \textit{Le Bitcoin est à la fois un protocole informatique qui utilise des technologies comme la cryptographie et la blockchain, et aussi une unité monétaire qui est utilisée pour les transactions.}

\item \textbf{Question 4:} A
\\ \textit{Ce cours sur Bitcoin se concentre sur les aspects monétaires de Bitcoin, y compris comment acheter et vendre des bitcoins, les stocker de manière sécurisée dans des portefeuilles numériques, et les utiliser pour des transactions.}

\item \textbf{Question 5:} B
\\ \textit{Bitcoin est une monnaie neutre, ce qui signifie qu'elle n'est contrôlée par aucune entité ou institution, donc chaque utilisateur a les mêmes droits. C'est un aspect clé de sa décentralisation.}

\item \textbf{Question 6:} D
\\ \textit{Les cypherpunks sont un groupe d'individus qui croyaient que la cryptographie pouvait servir d'outil pour protéger les droits individuels contre les ingérences des gouvernements et des grandes entreprises.}

\item \textbf{Question 7:} C
\\ \textit{Le Manifeste Cypherpunk a été écrit par Eric Hughes en 1993 et affirme que la vie privée est un droit fondamental.}

\item \textbf{Question 8:} A
\\ \textit{David Chaum a introduit le concept d'argent électronique anonyme avec son projet DigiCash dans les années 1980.}

\item \textbf{Question 9:} B
\\ \textit{La Déclaration de l'Indépendance du Cyberespace a été écrite par John Perry Barlow en 1996. Ce document affirme que le cyberespace est un domaine distinct de la sphère physique, donc les gouvernements industriels ne peuvent pas y imposer leur autorité.}

\item \textbf{Question 10:} C
\\ \textit{Le b-money de Wei Dai était l'idée d'une monnaie numérique anonyme sans entité centrale qui a grandement inspiré la création de Bitcoin. En fait, il est cité dans le livre blanc de Bitcoin.}

\item \textbf{Question 11:} B
\\ \textit{Au-delà de sa technologie, Bitcoin symbolise une révolution contre les systèmes financiers traditionnels : le protocole impose un nouveau système monétaire sans aucune entité centrale.}

\item \textbf{Question 12:} A
\\ \textit{Les premières formes de monnaie étaient souvent liées à des biens essentiels tels que les céréales, le bétail et d'autres marchandises. Cependant, ces biens présentaient des inconvénients majeurs, tels que la périssabilité, rendant difficile leur utilisation comme moyen d'épargne à long terme.}

\item \textbf{Question 13:} A
\\ \textit{Selon Aristote, la monnaie remplit trois fonctions : elle est une réserve de valeur, un moyen d'échange et une unité de compte. De nos jours, la fonction de moyen d'échange est largement reconnue comme la principale, avant les deux autres.}

\item \textbf{Question 14:} A
\\ \textit{Une monnaie efficace doit être fongible (interchangeable sans perte de valeur), divisible (séparable en petites parties pour faciliter les transactions de différents volumes), et liquide (facilement convertible en biens ou services).}

\item \textbf{Question 15:} B
\\ \textit{L'un des principaux inconvénients de l'or en tant que forme d'argent est qu'il est difficile de le diviser en petites unités selon les besoins qui se présentent, ce qui le rend moins pratique à utiliser.}

\item \textbf{Question 16:} D
\\ \textit{L'or est devenu le standard comme moyen d'échange en raison de sa rareté, de sa durabilité et de sa divisibilité : sa rareté empêchait la dévaluation causée par la production, sa durabilité assurait qu'il ne s'éroderait pas avec le temps, et sa divisibilité permettait son utilisation dans des transactions de différentes tailles.}

\item \textbf{Question 17:} B
\\ \textit{L'argent, en tant qu'outil de communication, permet la communication à travers l'espace et le temps, car nous transformons notre temps et notre énergie en un actif qui peut être réutilisé dans le futur sans risque de dévaluation. Il permet également la communication dans un langage commun universel, permettant à deux inconnus d'échanger, de commercer, et de s'accorder sur la valeur des choses.}

\item \textbf{Question 18:} B
\\ \textit{Le principal inconvénient des monnaies fiduciaires, telles que le dollar et l'euro, est que leur valeur est constamment érodée par l'inflation monétaire. Cela est dû à la dévaluation causée par les entités qui les contrôlent (roi, banque centrale, empereur, dictateur).}

\item \textbf{Question 19:} A
\\ \textit{Le principal avantage du Bitcoin par rapport aux autres formes de monnaie est qu'il offre une excellente réserve de valeur en raison de son offre strictement limitée et de son absence de tiers de confiance. Cela en fait un bon choix pour l'épargne à long terme.}

\item \textbf{Question 20:} C
\\ \textit{L'or ne peut pas suivre dans un monde globalisé et numérique parce qu'il nécessite une entité centrale pour le rendre divisible et facilement échangeable. Cela le rend moins pratique à utiliser dans un monde où les transactions sont de plus en plus numériques et décentralisées.}

\item \textbf{Question 21:} C
\\ \textit{Malgré ses propriétés, telles que son offre strictement limitée, le bitcoin n'est toujours pas largement accepté dans le commerce aujourd'hui, en raison du manque d'adoption par les individus, donc les commerçants ne sont pas encouragés à l'accepter.}

\item \textbf{Question 22:} D
\\ \textit{La monnaie est une technologie qui a évolué avec le temps, commençant par être représentée par des pierres brutes, se transformant en pièces métalliques, puis en billets de banque, cartes bancaires, et finalement en réseaux pair-à-pair avec l'utilisation de la blockchain Bitcoin et du Lightning Network.}

\item \textbf{Question 23:} D
\\ \textit{L'or est dense et lourd, ce qui signifie que transporter de grandes quantités peut être physiquement difficile. Ce poids le rend moins pratique pour le transport par rapport à des matériaux plus légers.}

\item \textbf{Question 24:} B
\\ \textit{Les monnaies fiduciaires, telles que le dollar et l'euro, sont très liquides et facilement transportables car elles sont maintenant majoritairement numériques. Cependant, leur valeur est constamment érodée par l'inflation monétaire.}

\item \textbf{Question 25:} B
\\ \textit{Le bitcoin, en raison de ses propriétés telles que sa divisibilité, sa vérifiabilité, son aspect sans permission et sa transportabilité, offre une excellente réserve de valeur. De plus, en tant que monnaie neutre d'internet, il représente un bon moyen d'échange qui ne connaît pas de frontières.}

\item \textbf{Question 26:} A
\\ \textit{La monnaie saine est un phénomène social naturel qui doit émerger du marché et du consensus volontaire. Aucun humain ou groupe d'humains ne peut créer de la monnaie.}

\item \textbf{Question 27:} A
\\ \textit{Bitcoin peut être vu comme une nouvelle forme d'argent, car il émerge d'un réseau pair-à-pair sans entité centrale.}

\item \textbf{Question 28:} C
\\ \textit{Bien qu'il ne soit pas facilement divisible ou transportable sur de longues distances, l'or est un matériau très durable, ce qui en fait une bonne réserve de valeur. Les pièces d'or étaient utilisées comme monnaie dans le passé, mais elles ne serviraient pas le même objectif à l'ère numérique actuelle.}

\item \textbf{Question 29:} B
\\ \textit{L'argent fiat est établi comme monnaie légale par décret gouvernemental, et sa valeur n'est pas dérivée d'une marchandise physique.}

\item \textbf{Question 30:} C
\\ \textit{Les monnaies fiduciaires ne sont pas soutenues par une marchandise physique comme l'or ou l'argent. Leur valeur est dérivée de la confiance que les gens ont dans les institutions qui les émettent et les régulent.}

\item \textbf{Question 31:} C
\\ \textit{Les banques centrales sont les entités responsables de l'émission d'une monnaie fiduciaire. Elles décident de la politique monétaire et déterminent la quantité d'argent qui doit être mise en circulation ou imprimée.}

\item \textbf{Question 32:} D
\\ \textit{L'offre totale de Bitcoin est plafonnée à 21 millions d'unités. C'est une caractéristique clé du Bitcoin qui le différencie des monnaies fiduciaires, qui peuvent être imprimées en quantités illimitées par les banques centrales.}

\item \textbf{Question 33:} D
\\ \textit{Le Bitcoin a été créé dans le but de séparer l'État de l'argent. C'est une réponse aux défis systémiques du système financier actuel, où l'État a un contrôle total sur l'offre monétaire et peut la manipuler à son propre avantage.}

\item \textbf{Question 34:} B
\\ \textit{Les entités centrales dévaluent souvent une monnaie en augmentant sa masse monétaire par rapport à la masse initiale de quelques pourcents chaque année. Cette dévaluation silencieuse est souvent justifiée comme étant dans l'intérêt du peuple.}

\item \textbf{Question 35:} C
\\ \textit{L'impression monétaire, ou la création de nouvelle monnaie par les banques centrales, conduit à l'inflation. Cela est dû au fait que l'augmentation de l'offre de monnaie réduit le pouvoir d'achat de la monnaie, entraînant une hausse des prix et une diminution de la valeur de l'épargne.}

\item \textbf{Question 36:} D
\\ \textit{Les monnaies numériques émises par les banques centrales (CBDCs) pourraient potentiellement entraver la liberté financière des individus et faciliter les abus autoritaires. Cela est dû au fait qu'elles offriraient une économie plus planifiée centralement, avec la banque centrale ayant encore plus de contrôle sur l'offre monétaire.}

\item \textbf{Question 37:} D
\\ \textit{Historiquement, les dirigeants ont centralisé l'or et introduit une nouvelle monnaie équivalente en valeur à l'or. Cette nouvelle monnaie est ensuite progressivement dévaluée, souvent sous le prétexte d'être dans l'intérêt du peuple.}

\item \textbf{Question 38:} B
\\ \textit{Une dévaluation soudaine ou une perte de confiance dans une monnaie fiduciaire peut conduire à l'hyperinflation. Il s'agit d'une situation où les prix des biens et services augmentent à un taux extrêmement élevé, menant souvent à une instabilité économique.}

\item \textbf{Question 39:} D
\\ \textit{En raison de sa nature décentralisée et cryptographique, Bitcoin ne peut pas être falsifié et est limité à 21 millions d'unités. C'est une caractéristique clé qui le différencie des monnaies fiduciaires, qui suivent une politique monétaire qui peut être modifiée par les banques centrales.}

\item \textbf{Question 40:} C
\\ \textit{Comme les acteurs politiques bénéficient du système financier actuel, ils ne sont pas incités à effectuer des changements radicaux. Cela permet au système de poursuivre sa trajectoire jusqu'à une possible implosion, car le système est régulé et restreint pour éviter son propre effondrement.}

\item \textbf{Question 41:} C
\\ \textit{L'hyperinflation est un phénomène monétaire caractérisé par une augmentation drastique de l'inflation à travers l'impression monétaire par les autorités, ce qui conduit à une perte complète de confiance dans une monnaie.}

\item \textbf{Question 42:} C
\\ \textit{À cause de l'hyperinflation, les individus perdent rapidement toute confiance dans la monnaie, ce qui conduit à une diminution drastique de sa valeur sur une courte période de temps.}

\item \textbf{Question 43:} A
\\ \textit{La première phase de l'hyperinflation est la perte de confiance dans une monnaie.}

\item \textbf{Question 44:} C
\\ \textit{La phase finale de l'hyperinflation se produit lorsqu'une nouvelle monnaie est introduite pour remplacer l'ancienne.}

\item \textbf{Question 45:} D
\\ \textit{Avec une inflation de 2\%, vous perdez 2\% de votre pouvoir d'achat annuellement, ce qui équivaut à 10\% sur 5 ans.}

\item \textbf{Question 46:} A
\\ \textit{Avec une inflation de 20\%, vous perdez presque la moitié de votre pouvoir d'achat en 3 ans.}

\item \textbf{Question 47:} C
\\ \textit{Un taux d'inflation de 100 \% par jour est un scénario réaliste qui s'est produit au Zimbabwe en 2007, où la valeur de l'argent doublait approximativement toutes les 24 heures.}

\item \textbf{Question 48:} A
\\ \textit{La deuxième phase de l'hyperinflation est l'effondrement de la monnaie & l'augmentation des prix, juste après la phase de Perte de Confiance.}

\item \textbf{Question 49:} C
\\ \textit{Le mois ayant connu le taux d'inflation le plus élevé en Allemagne était octobre 1923, atteignant un pic de 29 500 \%.}

\item \textbf{Question 50:} B
\\ \textit{Le mois ayant connu l'inflation la plus élevée en Hongrie était juillet 1946, avec une inflation record des prix de 41 900 000 000 000 000 \%.}

\item \textbf{Question 51:} B
\\ \textit{En novembre 2008, le taux d'inflation mensuel de la monnaie du Zimbabwe a atteint un pic de 79 600 000 000 \%, ce qui a réduit la valeur du dollar zimbabwéen à 0.}

\item \textbf{Question 52:} B
\\ \textit{En avril 2009, le Ministre des Finances a annoncé la suspension du dollar zimbabwéen et a autorisé l'utilisation de différentes devises étrangères pour le commerce.}

\item \textbf{Question 53:} C
\\ \textit{Le Bitcoin possède toutes les caractéristiques nécessaires pour être une monnaie efficace et saine : divisible, instantanément transportable, incensurable, coût de vérification négligeable, et avec une politique monétaire déjà établie pour les siècles à venir avec 21 millions d'unités.}

\item \textbf{Question 54:} B
\\ \textit{Le processus qui vérifie les transactions sur le réseau Bitcoin est appelé minage. Les mineurs effectuent une tâche cryptographique et sont récompensés par l'émission de nouveaux bitcoins.}

\item \textbf{Question 55:} B
\\ \textit{L'événement lors duquel la récompense de minage est divisée par deux est appelé halving. Cet événement donne à la courbe d'émission monétaire une forme d'escalier.}

\item \textbf{Question 56:} C
\\ \textit{Le mécanisme qui garantit qu'un nouveau bloc est ajouté à la blockchain toutes les dix minutes est l'ajustement de la difficulté. Cet ajustement a lieu tous les 2016 blocs, soit environ toutes les deux semaines.}

\item \textbf{Question 57:} D
\\ \textit{La théorie des jeux est un concept mathématique qui est basé sur la rationalité humaine. Dans Bitcoin, les changements de protocole sont appliqués par les utilisateurs, rendant toute modification du protocole Bitcoin extrêmement complexe.}

\item \textbf{Question 58:} D
\\ \textit{La commande pour vérifier la quantité de bitcoins en circulation sur un nœud Bitcoin est bitcoin-cli gettxoutsetinfo. Cette commande offre transparence et vérifiabilité, renforçant ainsi la confiance dans le système Bitcoin.}

\item \textbf{Question 59:} D
\\ \textit{Bitcoin s'aligne sur les principes de l'économie autrichienne. Sa quantité limitée et l'absence d'entité centrale le protègent des risques inhérents d'inflation que l'on observe dans les monnaies traditionnelles.}

\item \textbf{Question 60:} B
\\ \textit{En raison du mécanisme de halving, il est mathématiquement prévisible que la création de bitcoins cessera en 2140, lorsque le nombre total de bitcoins atteindra sa limite de 21 millions.}

\item \textbf{Question 61:} B
\\ \textit{L'ajustement de la difficulté est un mécanisme qui a lieu tous les 2016 blocs, soit environ toutes les deux semaines, pour garantir qu'un nouveau bloc soit ajouté à la blockchain toutes les dix minutes en moyenne, indépendamment du hashrate global.}

\item \textbf{Question 62:} B
\\ \textit{La récompense pour le minage est programmée pour être divisée par deux tous les 210 000 blocs, ce qui correspond à peu près à tous les quatre ans : un événement connu sous le nom de halving.}

\item \textbf{Question 63:} C
\\ \textit{Après le premier halving, qui a eu lieu à la hauteur de bloc 210 000, le nombre estimé de bitcoins en circulation était de 10 500 000 BTC, ce qui représente la moitié du nombre total de bitcoins qui seront en circulation.}

\item \textbf{Question 64:} D
\\ \textit{Après le quatrième halving, qui se produit à la hauteur de bloc 840 000, la récompense de minage est de 3.125 BTC.}

\item \textbf{Question 65:} A
\\ \textit{Le 7ème halving devrait avoir lieu en 2036, et le nombre estimé de bitcoins en circulation sera de 20,835,937.5, ce qui correspondra approximativement à 99.22\% de l'offre totale.}

\item \textbf{Question 66:} A
\\ \textit{Un portefeuille Bitcoin est utilisé pour interagir avec le réseau Bitcoin. Ses trois principales fonctions sont de permettre de recevoir des bitcoins, d'envoyer des bitcoins, et de les sécuriser contre les tentatives de piratage et de vol.}

\item \textbf{Question 67:} D
\\ \textit{Lors de l'initialisation d'un portefeuille Bitcoin, une liste secrète de mots, appelée graine mnémonique, est générée. Cette liste de récupération est très importante car elle représente la propriété des Bitcoins et donc le droit de les envoyer.}

\item \textbf{Question 68:} A
\\ \textit{Bien que vos clés privées soient stockées dans votre portefeuille, les bitcoins eux-mêmes sont en réalité stockés dans la blockchain Bitcoin, qui est le registre public distribué au sein du réseau pair-à-pair Bitcoin.}

\item \textbf{Question 69:} A
\\ \textit{Depuis 2017, la clé privée peut être codée dans une simple liste de 12 ou 24 mots, appelée la phrase mnémonique. Cette phrase est une sauvegarde de votre portefeuille Bitcoin, donc elle vous permet d'accéder à vos bitcoins en utilisant n'importe quel logiciel/application de portefeuille Bitcoin.}

\item \textbf{Question 70:} D
\\ \textit{La clé publique est dérivée de la clé privée et est utilisée pour générer des adresses Bitcoin afin de recevoir de l'argent.}

\item \textbf{Question 71:} A
\\ \textit{Partager la clé publique comporte des risques pour la confidentialité, mais pas pour la sécurité. Cela est dû au fait que la clé publique est utilisée pour générer des adresses Bitcoin et donc recevoir de l'argent, mais elle ne représente pas la propriété des bitcoins.}

\item \textbf{Question 72:} A
\\ \textit{Perdre l'appareil sur lequel se trouve votre portefeuille ne signifie pas nécessairement la perte de vos bitcoins. Ce qui vous permet de recréer votre portefeuille et de dépenser vos bitcoins est la phrase mnémonique.}

\item \textbf{Question 73:} C
\\ \textit{Si vous n'êtes pas satisfait de la sécurité par défaut de votre portefeuille Bitcoin, vous pouvez toujours la renforcer en ajoutant une phrase secrète à votre portefeuille Bitcoin.}

\item \textbf{Question 74:} D
\\ \textit{Grâce à la cryptographie utilisée pour créer le portefeuille, la probabilité que quelqu'un devine accidentellement votre liste de 12 ou 24 mots est extrêmement faible. C'est comme trouver le bon numéro entre 1 et $2^256$, ce qui est presque équivalent à trouver un atome défini dans l'Univers.}

\item \textbf{Question 75:} A
\\ \textit{Lors du choix d'un portefeuille Bitcoin, la question centrale est : êtes-vous le propriétaire des fonds ou laissez-vous le contrôle de votre argent à un tiers ? Cela est dû au fait que le détenteur de la clé privée est le propriétaire des bitcoins.}

\item \textbf{Question 76:} B
\\ \textit{Pour dépenser des bitcoins, le propriétaire doit utiliser l'Algorithme de Signature Numérique à Courbe Elliptique (ECDSA) pour forger une signature permettant la propagation d'une transaction Bitcoin.}

\item \textbf{Question 77:} A
\\ \textit{L'analogie utilisée pour expliquer comment vous pouvez recevoir des bitcoins sans permettre à l'expéditeur de voler vos fonds est une boîte aux lettres. Les gens déposent de l'argent dedans, mais vous êtes le seul à pouvoir l'ouvrir.}

\item \textbf{Question 78:} D
\\ \textit{Lorsque vous possédez des bitcoins, la préoccupation principale est la sécurité de vos fonds. Cela est dû au fait que le Bitcoin, comme toute autre forme de monnaie, peut être volé ou perdu s'il n'est pas correctement sécurisé.}

\item \textbf{Question 79:} D
\\ \textit{Un portefeuille chaud est un terme utilisé pour décrire un portefeuille Bitcoin qui est stocké sur un appareil ayant accès à internet, tel qu'un téléphone ou un ordinateur. De cette manière, il est facile de réaliser des transactions, mais cela signifie également que le portefeuille est potentiellement vulnérable aux menaces en ligne.}

\item \textbf{Question 80:} A
\\ \textit{Un portefeuille froid est un terme utilisé pour décrire un portefeuille Bitcoin qui est stocké sur un dispositif qui n'est pas connecté à internet. Cela le rend moins vulnérable aux menaces en ligne, mais cela signifie également que le portefeuille est moins accessible pour les transactions quotidiennes.}

\item \textbf{Question 81:} C
\\ \textit{Un portefeuille multisig est un type de portefeuille Bitcoin qui requiert plusieurs signatures pour réaliser une transaction. Cela ajoute une couche supplémentaire de sécurité, car cela signifie qu'une seule personne ne peut pas effectuer une transaction sans l'approbation des autres.}

\item \textbf{Question 82:} D
\\ \textit{Lorsque vous utilisez un service de garde pour stocker vos bitcoins, vous n'êtes pas le seul détenteur de vos bitcoins. Cela signifie que le tiers de confiance peut restreindre votre accès à vos fonds à tout moment, vous donnant le même niveau de souveraineté financière que le système bancaire traditionnel.}

\item \textbf{Question 83:} A
\\ \textit{Compliquer excessivement la sécurité et l'accessibilité de vos bitcoins peut vous nuire si, par exemple, vous manipulez mal les sauvegardes de vos portefeuilles. Cela est dû au fait que des procédures de sauvegarde simples sont cruciales pour récupérer l'accès à vos fonds en cas de perte de votre téléphone ou ordinateur.}

\item \textbf{Question 84:} A
\\ \textit{Un portefeuille mobile est recommandé pour les dépenses quotidiennes car il est facile de réaliser des transactions avec. Cependant, il n'est pas recommandé pour stocker de grandes sommes ou pour les économies à long terme, car il est potentiellement vulnérable aux menaces en ligne.}

\item \textbf{Question 85:} D
\\ \textit{Un portefeuille physique hors ligne, ou cold wallet, est recommandé pour stocker de grandes quantités car il est moins vulnérable aux menaces en ligne que les portefeuilles chauds (hot wallet). Cependant, il n'est pas recommandé pour les dépenses quotidiennes car il est moins accessible pour les transactions fréquentes.}

\item \textbf{Question 86:} D
\\ \textit{Un portefeuille multisig est recommandé pour stocker de grandes quantités et pour un usage corporatif car il nécessite plusieurs signatures pour effectuer une transaction. Ce processus ajoute une couche supplémentaire de sécurité, car une seule personne ne peut pas réaliser une transaction sans l'approbation des autres.}

\item \textbf{Question 87:} C
\\ \textit{Lorsque vous utilisez un portefeuille avec une phrase secrète, vous devez sauvegarder à la fois la liste de 12 ou 24 mots et votre phrase secrète, car la liste de mots et la phrase secrète sont tous les deux nécessaires pour récupérer l'accès à vos fonds en cas de perte de votre téléphone ou ordinateur. Sans l'un ou l'autre, les fonds ne sont pas accessibles.}

\item \textbf{Question 88:} D
\\ \textit{Lorsque vous utilisez un portefeuille multisig, chaque clé privée du multisig devrait être stockée dans des emplacements différents. Cela est dû au fait que le multisig nécessite plusieurs signatures pour réaliser une transaction, donc un certain nombre de portefeuilles sont nécessaires. Ainsi, stocker chaque partie dans un emplacement différent ajoute une couche supplémentaire de sécurité.}

\item \textbf{Question 89:} A
\\ \textit{Le terme hodl provient d'une publication en ligne mal orthographiée en 2013 sur un forum Bitcoin. Avec le temps, hodl a évolué pour représenter une stratégie d'investissement plus large. Il encourage les investisseurs à conserver leurs cryptomonnaies sur le long terme, indépendamment des fluctuations du marché.}

\item \textbf{Question 90:} B
\\ \textit{Dans un portefeuille Bitcoin, la clé privée est souvent représentée par une liste de 12 ou 24 mots, également connue sous le nom de phrase de récupération ou phrase mnémonique.}

\item \textbf{Question 91:} A
\\ \textit{Si votre clé privée est révélée à un tiers, les fonds associés ne sont plus sécurisés, car la clé privée leur permet d'accéder à vos fonds.}

\item \textbf{Question 92:} D
\\ \textit{Une des alternatives recommandées au papier pour stocker votre phrase mnémonique est de la graver sur une plaque métallique. Cette solution est plus durable et résistante dans le temps.}

\item \textbf{Question 93:} D
\\ \textit{Une fois les mots de votre phrase mnémonique écrits, il est recommandé de faire une seconde copie et de la stocker dans un endroit séparé de la première. Cela fournit une sauvegarde en cas de perte ou d'accident avec la première copie.}

\item \textbf{Question 94:} D
\\ \textit{L'absence d'une liste de 12 ou 24 mots, c'est-à-dire la phrase mnémonique, dans un portefeuille devrait vous alerter que le portefeuille n'est pas sécurisé car il ne suit pas les procédures de sauvegarde standard.}

\item \textbf{Question 95:} A
\\ \textit{La méthode standard pour récupérer la clé privée d'un portefeuille Bitcoin consiste à saisir votre phrase mnémonique dans n'importe quel logiciel de portefeuille. Cela vous permet de restaurer votre portefeuille.}

\item \textbf{Question 96:} B
\\ \textit{Graver votre phrase mnémonique sur un matériau résistant comme l'acier crée une sauvegarde physique de vos clés qui est résistante à la fois aux dégâts des eaux et au feu, sécurisant ainsi vos bitcoins sur le long terme.}

\item \textbf{Question 97:} C
\\ \textit{Un plan de succession garantit que vos bitcoins seront correctement gérés après votre décès. Il peut inclure une lettre manuscrite détaillant vos actifs, leurs méthodes d'accès, et les coordonnées de personnes de confiance à contacter.}

\item \textbf{Question 98:} B
\\ \textit{Une phrase mnémonique dans le contexte du Bitcoin est une liste de 12 ou 24 mots qui vous permet de restaurer votre portefeuille sur n'importe quelle application de portefeuille Bitcoin. Quiconque ayant accès à cette liste peut également accéder à vos bitcoins.}

\item \textbf{Question 99:} D
\\ \textit{La confidentialité est importante dans le contexte du Bitcoin pour minimiser les risques que vos fonds soient tracés, ce qui peut faciliter votre surveillance.}

\item \textbf{Question 100:} A
\\ \textit{Il est conseillé d'éviter de laisser vos bitcoins sur des plateformes d'échange car ils peuvent être susceptibles aux attaques de hackers.}

\item \textbf{Question 101:} A
\\ \textit{Dans le contexte de l'héritage de Bitcoin, il est important de discuter de l'héritage des bitcoins avec un notaire pour assurer la conformité fiscale.}

\item \textbf{Question 102:} C
\\ \textit{Les portefeuilles Bitcoin sont un type de logiciel utilisé pour stocker les clés privées qui vous permettent de réaliser vos transactions en bitcoins.}

\item \textbf{Question 103:} B
\\ \textit{Dans Bitcoin, la souveraineté financière va de pair avec la responsabilité individuelle, car vous devez prendre en charge la garde de vos fonds en suivant les procédures de sauvegarde recommandées.}

\item \textbf{Question 104:} B
\\ \textit{Le Blockmit est une solution peu coûteuse pour graver votre phrase mnémonique sur un matériau résistant comme l'acier, sécurisant ainsi vos bitcoins sur le long terme.}

\item \textbf{Question 105:} D
\\ \textit{KYC (Know Your Customer) ou Connaître Votre Client fait référence à la procédure que les entreprises adoptent pour vérifier l'identité de leurs clients. Cela est généralement fait pour assurer la conformité avec les réglementations et pour prévenir des activités illégales telles que le blanchiment d'argent ou la fraude.}

\item \textbf{Question 106:} A
\\ \textit{Acheter des bitcoins non-KYC signifie acquérir des bitcoins sans passer par les procédures de vérification standard Connaître Votre Client (Know Your Customer) que de nombreuses bourses et plateformes exigent. Une raison principale pour laquelle les individus pourraient choisir cette voie est de maintenir leur vie privée, en s'assurant que leurs informations personnelles et activités financières restent non divulgués.}

\item \textbf{Question 107:} A
\\ \textit{Alors qu'une personne pourrait prioriser la sécurité, une autre pourrait valoriser la convivialité ou des fonctionnalités spécifiques. Par conséquent, un seul portefeuille Bitcoin pourrait ne pas répondre aux besoins de tous, rendant nécessaire l'existence de multiples options de portefeuilles.}

\item \textbf{Question 108:} B
\\ \textit{Le Bitcoin a été présenté au monde pour la première fois le 31 octobre 2008, lorsque Satoshi Nakamoto a envoyé un courriel contenant le Livre Blanc de Bitcoin à la liste de diffusion cypherpunk.}

\item \textbf{Question 109:} A
\\ \textit{Hal Finney était la première personne à rejoindre le réseau Bitcoin après Satoshi Nakamoto lui-même. Il a même tweeté "Running bitcoin" pour marquer l'occasion.}

\item \textbf{Question 110:} B
\\ \textit{La première transaction Bitcoin enregistrée était une transaction de 10 BTC entre Satoshi Nakamoto et Hal Finney.}

\item \textbf{Question 111:} C
\\ \textit{Satoshi Nakamoto a disparu en 2010, laissant sa création entre les mains de la communauté.}

\item \textbf{Question 112:} A
\\ \textit{Le premier bloc de la blockchain Bitcoin, également connu sous le nom de bloc Genesis, contient le message "03/jan/2009 Le Chancelier au bord d'un second sauvetage des banques", qui est le titre du *Times* ce jour-là pour horodater le bloc Genesis.}

\item \textbf{Question 113:} D
\\ \textit{Le forum BitcoinTalk est un lieu de discussion pour les utilisateurs de Bitcoin. Il a été créé par Satoshi Nakamoto le 22 novembre 2009.}

\item \textbf{Question 114:} C
\\ \textit{La première fois que le Bitcoin a été utilisé pour acheter des biens physiques fut quand Laszlo Hanyecz a acheté 2 pizzas pour 10 000 BTC, le 22 mai 2010, date qui est maintenant connue sous le nom de "Pizza Day".}

\item \textbf{Question 115:} A
\\ \textit{La disponibilité en ligne de Bitcoin est de 99,988 \% depuis 2009.}

\item \textbf{Question 116:} A
\\ \textit{La disponibilité en ligne de Bitcoin est de 99,988 \% depuis 2009.}

\item \textbf{Question 117:} A
\\ \textit{La première transaction Bitcoin, qui était une transaction de 10 BTC entre Satoshi Nakamoto et Hal Finney, est enregistrée dans le bloc 170.}

\item \textbf{Question 118:} C
\\ \textit{Le premier logiciel Bitcoin publié était Bitcoin-0.1.0. Il a été annoncé par Satoshi Nakamoto le 8 janvier 2009.}

\item \textbf{Question 119:} D
\\ \textit{La date du dernier message privé connu de Satoshi Nakamoto par email est le 23 avril 2011. Il a été envoyé au développeur Mike Hearn.}

\item \textbf{Question 120:} C
\\ \textit{Une transaction Bitcoin est un transfert numérique de la propriété de bitcoins, grâce à l'utilisation d'une adresse Bitcoin dérivée grâce à la cryptographie de Courbe Elliptique. Cette adresse est unique au portefeuille de la personne recevant les bitcoins.}

\item \textbf{Question 121:} C
\\ \textit{Les frais de transaction dans les transactions Bitcoin servent d'incitation pour les mineurs à inclure la transaction dans le prochain bloc. Les frais créent un marché libre pour l'inclusion des transactions dans les blocs.}

\item \textbf{Question 122:} D
\\ \textit{Créer un plan de succession est une étape cruciale pour assurer que vos bitcoins soient correctement gérés après votre décès. Ce plan peut inclure une lettre manuscrite listant vos actifs, leurs méthodes d'accès, et les informations des personnes de confiance à contacter.}

\item \textbf{Question 123:} A
\\ \textit{La blockchain Bitcoin est un registre public et immuable de toutes les transactions Bitcoin. C'est un registre commun pour tous les utilisateurs de Bitcoin.}

\item \textbf{Question 124:} A
\\ \textit{Le nombre de confirmations dans une transaction Bitcoin indique combien de blocs ont été minés au-dessus du bloc contenant la transaction. Plus une transaction a de confirmations, plus elle devient immuable.}

\item \textbf{Question 125:} D
\\ \textit{Un nœud Bitcoin dans le réseau Bitcoin s'occupe de la validation des nouveaux blocs et des transactions incluses. C'est ce réseau de nœuds Bitcoin qui le rend véritablement décentralisé.}

\item \textbf{Question 126:} B
\\ \textit{La répartition géographique des nœuds Bitcoin rend pratiquement impossible de détruire toutes les copies de la blockchain. Étant donné que les nœuds sont dispersés, il est difficile de tous les saisir physiquement.}

\item \textbf{Question 127:} B
\\ \textit{Les mineurs valident les transactions et les incluent dans un bloc à travers un processus appelé preuve de travail (proof of work). Cela implique de résoudre une énigme cryptographique.}

\item \textbf{Question 128:} C
\\ \textit{Dans le contexte du minage de blocs dans le réseau Bitcoin, un hash valide est une empreinte digitale unique associée au bloc, composée de 256 caractères. La validité de ce hash dépend de la difficulté du réseau Bitcoin.}

\item \textbf{Question 129:} A
\\ \textit{L'ajustement de difficulté dans le réseau Bitcoin assure qu'un bloc est ajouté environ toutes les dix minutes. Cela fait partie des règles du protocole du réseau Bitcoin.}

\item \textbf{Question 130:} B
\\ \textit{Le réseau Bitcoin s'assure qu'il est trop coûteux d'agir de manière malveillante par des règles de protocole et des incitations économiques. Cela élimine le besoin de confiance à n'importe quelle étape du transfert de propriété du bitcoin.}

\item \textbf{Question 131:} B
\\ \textit{Les utilisateurs dans le réseau Bitcoin transfèrent la propriété de leur argent en signant numériquement les transactions avec leurs clés privées. Cela vérifie qu'ils sont les propriétaires des bitcoins qu'ils souhaitent transférer.}

\item \textbf{Question 132:} A
\\ \textit{Les nœuds dans le réseau Bitcoin remplissent plusieurs fonctions cruciales, y compris maintenir une copie de la blockchain Bitcoin, valider les transactions, transmettre les informations aux autres nœuds et faire respecter les règles du protocole Bitcoin.}

\item \textbf{Question 133:} A
\\ \textit{En septembre 2023, il y a environ 45 000 nœuds répartis à travers le monde.}

\item \textbf{Question 134:} C
\\ \textit{La blockchain Bitcoin est d'environ 660Go en 2025.}

\item \textbf{Question 135:} D
\\ \textit{Faire fonctionner un nœud Bitcoin sur un Raspberry Pi 4 coûte un peu moins de 10€ par an en termes de consommation électrique en 2023.}

\item \textbf{Question 136:} B
\\ \textit{Un nœud élagué dans le contexte du Bitcoin est un nœud qui ne conserve que les derniers N blocs en mémoire.}

\item \textbf{Question 137:} D
\\ \textit{Si un utilisateur de nœud décidait de suivre d'autres règles qui invalident les règles de consensus actuelles, alors ce nœud ne ferait plus partie du réseau Bitcoin.}

\item \textbf{Question 138:} A
\\ \textit{En considérant 1 bloc de 1MB toutes les 10 minutes, exécuter un nœud Bitcoin nécessite environ 5GB par mois en bande passante.}

\item \textbf{Question 139:} C
\\ \textit{La guerre des blocs en 2017 était un conflit sur la question de savoir s'il fallait modifier Bitcoin en augmentant la taille des blocs pour étendre la capacité de transaction.}

\item \textbf{Question 140:} A
\\ \textit{Le coût abordable et l'accessibilité d'un nœud Bitcoin sont importants car ils facilitent la décentralisation du réseau.}

\item \textbf{Question 141:} C
\\ \textit{Si les blocs étaient 100 fois plus lourds, faire fonctionner un nœud Bitcoin nécessiterait un disque dur de 70TB, une bande passante de plus de 500GB/mois, et un matériel capable de valider des centaines de milliers de transactions en moins de 10 minutes. Cela rendrait l'exploitation d'un nœud Bitcoin inaccessible à la personne moyenne, compromettant la décentralisation du protocole et l'immuabilité des transactions et des règles de consensus.}

\item \textbf{Question 142:} A
\\ \textit{Lors de la guerre des blocs en 2017, les utilisateurs et les nœuds ont triomphé en rejetant la proposition de modification visant à augmenter la taille des blocs, et ont activé une mise à jour appelée SegWit.}

\item \textbf{Question 143:} B
\\ \textit{Le Lightning Network est un réseau de paiement Bitcoin instantané qui est basé sur la blockchain Bitcoin.}

\item \textbf{Question 144:} A
\\ \textit{Les mineurs utilisent leur puissance de calcul pour résoudre des problèmes mathématiques complexes, ce qui leur permet d'ajouter de nouvelles transactions à la blockchain. Ce processus sécurise le réseau et maintient son intégrité.}

\item \textbf{Question 145:} B
\\ \textit{La Preuve de Travail est un algorithme de consensus utilisé dans le réseau Bitcoin pour valider les transactions et créer de nouveaux blocs. Elle exige des mineurs de résoudre des problèmes mathématiques complexes, assurant ainsi la sécurité et l'intégrité du réseau.}

\item \textbf{Question 146:} C
\\ \textit{Le HashRate fait référence au nombre total de tentatives que les mineurs font par seconde pour résoudre la Preuve de Travail. C'est une mesure de la puissance de calcul du réseau Bitcoin.}

\item \textbf{Question 147:} D
\\ \textit{L'ajustement de difficulté dans le minage de Bitcoin est conçu pour garantir que le temps nécessaire pour miner un bloc reste approximativement constant, indépendamment de la puissance de calcul totale du réseau. Il ajuste la difficulté du problème mathématique que les mineurs doivent résoudre, en fonction du nombre de mineurs participant au réseau.}

\item \textbf{Question 148:} A
\\ \textit{Les ASICs (Circuits Intégrés Spécifiques à une Application) sont du matériel spécialisé conçu spécifiquement pour le minage de Bitcoin. Ils sont optimisés pour effectuer la fonction de hachage SHA256, qui est le problème mathématique que les mineurs doivent résoudre dans le processus de Preuve de Travail (Proof of Work).}

\item \textbf{Question 149:} D
\\ \textit{La transaction coinbase est toujours la première transaction d'un bloc. Elle inclut la récompense du mineur pour avoir effectué le travail de validateur.}

\item \textbf{Question 150:} C
\\ \textit{Les pools de minage sont des groupes de mineurs qui combinent leur puissance de calcul pour augmenter leurs chances de miner un bloc. En regroupant leurs ressources, ils peuvent miner des blocs plus fréquemment et recevoir des récompenses plus petites mais plus régulières.}

\item \textbf{Question 151:} B
\\ \textit{Le problème des généraux byzantins est un problème classique en informatique distribuée, qui implique d'atteindre un consensus dans un réseau où certains des nœuds peuvent être défectueux ou malveillants. Bitcoin résout ce problème grâce à l'utilisation d'une blockchain basée sur la Preuve de Travail, qui assure que tous les nœuds sont d'accord sur l'état de l'historique des transactions.}

\item \textbf{Question 152:} D
\\ \textit{Le Proof of Work est un algorithme de consensus utilisé dans le réseau Bitcoin pour valider les transactions et créer de nouveaux blocs. Il exige des mineurs de résoudre des problèmes mathématiques complexes, ce qui assure l'authenticité des informations dans le réseau sans la nécessité d'un tiers de confiance.}

\item \textbf{Question 153:} D
\\ \textit{Le coût énergétique dans le minage de Bitcoin est un aspect crucial de la sécurité du réseau. Le coût énergétique élevé du minage rend extrêmement coûteuse toute tentative de manipulation de la blockchain par un acteur malveillant, tandis que le faible coût de vérification des transactions assure que n'importe qui peut participer au réseau.}

\item \textbf{Question 154:} D
\\ \textit{Dans une attaque à 51\%, un agent avec plus de la moitié du hashrate du réseau pourrait potentiellement manipuler la blockchain en créant une chaîne alternative de blocs. Cependant, une telle attaque serait extrêmement coûteuse et peu probable d'être rentable, rendant ce scénario peu probable.}

\item \textbf{Question 155:} B
\\ \textit{La principale préoccupation environnementale du minage de Bitcoin est la consommation d'énergie. Bien que les ASICs soient principalement fabriqués en aluminium et puissent être recyclés ou réutilisés, l'énergie nécessaire pour alimenter ces machines est significative.}

\item \textbf{Question 156:} C
\\ \textit{Les mineurs dans le protocole Bitcoin sécurisent à la fois le réseau actuel et historique. Ils ne régulent pas le prix du Bitcoin, ne contrôlent pas son offre, ni ne gèrent les transactions entre utilisateurs.}

\item \textbf{Question 157:} D
\\ \textit{En moyenne, un bloc est créé toutes les 10 minutes dans le protocole Bitcoin. Ce temps reste constant indépendamment du nombre de mineurs et de leur puissance de calcul.}

\item \textbf{Question 158:} B
\\ \textit{Bitcoin offre une forme de liberté financière en échappant à la censure et aux restrictions bancaires. Cela est particulièrement bénéfique pour les individus vivant sous une oppression financière ou un régime dictatorial.}

\item \textbf{Question 159:} B
\\ \textit{Le protocole Bitcoin assure que le temps moyen entre deux blocs reste constant en ajustant la difficulté de trouver un hash valide tous les 2016 blocs. Cet événement se produit indépendamment du nombre de mineurs et de leur puissance de calcul.}

\item \textbf{Question 160:} C
\\ \textit{Les mineurs peuvent agir comme acheteur de dernier recours pour les centrales électriques qui ne sont pas encore connectées au réseau. Cela permet aux centrales électriques de sécuriser un financement même avant d'être connectées au réseau électrique.}

\item \textbf{Question 161:} B
\\ \textit{Le système financier actuel est critiqué pour encourager la surconsommation et l'endettement, ce qui pose de sérieux problèmes environnementaux. Ceci est dû à l'accès facile au crédit, à l'émission monétaire par les banques, et à l'utilisation du système de réserve fractionnaire.}

\item \textbf{Question 162:} A
\\ \textit{Le Bitcoin favorise l'utilisation d'énergies vertes, ce qui peut potentiellement aider à la transition écologique. Par exemple, les flammes dans les puits de pétrole, qui brûlent le méthane pour diminuer la pollution, peuvent être éteintes par les mineurs de Bitcoin.}

\item \textbf{Question 163:} D
\\ \textit{Le protocole Bitcoin récompense le mineur qui trouve un hash valide pour le bloc suivant avec une récompense dont le montant est défini par les règles de consensus, en plus des frais de toutes les transactions incluses dans le bloc valide.}

\item \textbf{Question 164:} A
\\ \textit{Les machines spécialisées SHA256, ou ASICs, augmentent le taux de hachage par joule dans le minage de Bitcoin. Cela signifie le nombre de tentatives par seconde et par énergie consommée.}

\item \textbf{Question 165:} A
\\ \textit{Le protocole Bitcoin assure sa résistance à la censure et sa résilience parce que chaque composant du protocole est distribué géographiquement à travers le monde. Par exemple, il y a environ 45 000 nœuds Bitcoin sur l'ensemble des continents.}

\item \textbf{Question 166:} D
\\ \textit{La principale différence entre les idées de l'économie keynésienne et de l'économie autrichienne en relation avec le système monétaire est que l'économie keynésienne ne prend pas en compte les aspects temporels et dynamiques des situations et des ressources. En d'autres termes, une monnaie illimitée ne peut pas refléter efficacement les ressources limitées de notre planète.}

\item \textbf{Question 167:} D
\\ \textit{La nature décentralisée du Bitcoin, son pseudonymat et sa résistance à la censure le rendent attrayant pour des groupes tels que les technophiles, les cypherpunks, les libertariens, et les enthousiastes de l'or.}

\item \textbf{Question 168:} C
\\ \textit{Miner du Bitcoin nécessite une puissance de calcul substantielle, ce qui, à son tour, consomme de l'électricité. De plus, les mineurs doivent investir dans du matériel spécialisé pour miner efficacement.}

\item \textbf{Question 169:} B
\\ \textit{La volatilité du Bitcoin peut être atténuée par plusieurs solutions telles que les couvertures financières (stablecoins), une forte croyance à long terme (hodling), ou simplement ne pas mettre 100\% de son argent dans le Bitcoin sans rien comprendre.}

\item \textbf{Question 170:} D
\\ \textit{Le protocole Bitcoin est unique en ce qu'il fonctionne à l'échelle mondiale, 24 heures sur 24, 7 jours sur 7, rendant la régulation difficile pour les autorités financières.}

\item \textbf{Question 171:} B
\\ \textit{L'année 2017 a été marquée par une bulle spéculative significative dans le monde des cryptomonnaies, notamment avec le lancement de milliers d'Offres Initiales de Pièces (ICO).}

\item \textbf{Question 172:} B
\\ \textit{Le Bitcoin, en raison de sa quantité limitée à 21 millions de bitcoins et de son processus de halving (réduisant de moitié la création monétaire tous les 4 ans en moyenne), suit une tendance générale à la hausse.}

\item \textbf{Question 173:} B
\\ \textit{Le Bitcoin connaît de nombreuses bulles spéculatives parce que c'est une monnaie qui est réellement utilisée partout dans le monde.}

\item \textbf{Question 174:} D
\\ \textit{Le protocole Bitcoin est unique en ce qu'il fonctionne à l'échelle mondiale, 24 heures sur 24, 7 jours sur 7, rendant la régulation difficile pour les autorités financières.}

\item \textbf{Question 175:} B
\\ \textit{Le Bitcoin continue de survivre et même de se développer davantage en s'intégrant de plus en plus au marché traditionnel. L'arrivée des ETF Bitcoin, une réglementation plus claire, et une meilleure acquisition d'outils de stockage ne font qu'encourager cette tendance.}

\item \textbf{Question 176:} A
\\ \textit{L'utilisation de Bitcoin s'est étendue aux marchés noirs du web comme Silk Road en raison de sa nature incontrôlable et pseudonyme.}

\item \textbf{Question 177:} C
\\ \textit{Le prix du Bitcoin est souvent caractérisé par une volatilité significative, qui peut être influencée par les variations du marché et les phases haussières et baissières du marché.}

\item \textbf{Question 178:} A
\\ \textit{Le Bitcoin n'est pas une bulle qui va disparaître mais une monnaie qui est réellement utilisée partout dans le monde.}

\item \textbf{Question 179:} B
\\ \textit{Le Bitcoin est considéré comme une économie parallèle aux monnaies fiduciaires, ce qui signifie qu'il fonctionne à côté des économies traditionnelles et peut être utilisé pour des transactions directement, sans avoir besoin de passer par une plateforme d'échange.}

\item \textbf{Question 180:} C
\\ \textit{Une manière d'obtenir des bitcoins sans les acheter est de les accepter comme moyen de paiement pour les produits ou services que vous proposez. Cela vous permet d'acquérir des bitcoins à travers votre travail.}

\item \textbf{Question 181:} A
\\ \textit{Accepter le Bitcoin en tant que commerçant présente plusieurs avantages, y compris la résistance à la censure, la réduction des frais de transaction, l'augmentation de l'efficacité, la protection contre l'inflation, ainsi que la liberté financière et la souveraineté.}

\item \textbf{Question 182:} B
\\ \textit{BTCMap est un projet open-source et collaboratif qui répertorie tous les commerçants qui acceptent le Bitcoin ainsi que les différentes communautés Bitcoin autour du monde.}

\item \textbf{Question 183:} A
\\ \textit{Le problème des généraux byzantins est un problème classique dans le domaine du calcul distribué, qui implique d'atteindre un consensus dans un réseau où certains des nœuds peuvent être défectueux ou malveillants. Bitcoin résout ce problème par l'utilisation d'une blockchain basée sur la Preuve de Travail, ce qui assure que tous les nœuds s'accordent sur l'état de l'historique des transactions.}

\item \textbf{Question 184:} A
\\ \textit{L'AMF, ou Autorité des Marchés Financiers, est une organisation française qui régule les plateformes où les bitcoins peuvent être achetés.}

\item \textbf{Question 185:} B
\\ \textit{Lors du choix d'une solution pour accepter le Bitcoin pour votre entreprise, plusieurs facteurs doivent être pris en compte, tels que le volume de transactions attendu, le budget alloué, et le type d'entreprise (en ligne ou physique).}

\item \textbf{Question 186:} A
\\ \textit{OpenNode est une solution en ligne simple pour que les entreprises acceptent le Bitcoin. Parmi les autres solutions, on trouve Swiss Bitcoin Pay pour les commerçants amateurs et BTCpay Server pour les grandes structures ou les passionnés de Bitcoin.}

\item \textbf{Question 187:} D
\\ \textit{BTCMap est l'une des initiatives qui contribuent à rendre l'économie Bitcoin plus accessible et pratique pour tous. C'est un projet open-source et collaboratif qui répertorie tous les commerçants acceptant le Bitcoin ainsi que les différentes communautés Bitcoin à travers le monde.}

\item \textbf{Question 188:} C
\\ \textit{BTCpay Server est un outil permettant d'accepter le Bitcoin pour les grandes structures ou les passionnés de Bitcoin. D'autres solutions incluent OpenNode pour les entreprises en ligne simples et Swiss Bitcoin Pay pour les commerçants amateurs.}

\item \textbf{Question 189:} C
\\ \textit{Une manière de faciliter l'utilisation du Bitcoin dans les transactions quotidiennes est de lister tous les commerçants qui acceptent Bitcoin. Cela peut être réalisé à travers des plateformes comme BTCMap.}

\item \textbf{Question 190:} A
\\ \textit{Swiss Bitcoin Pay est une solution permettant aux commerçants amateurs d'accepter le Bitcoin. D'autres solutions incluent OpenNode pour les entreprises en ligne simples et BTCpay Server pour les grandes structures ou les passionnés de Bitcoin.}

\item \textbf{Question 191:} A
\\ \textit{Le Bitcoin est un actif financier hautement volatil, et son prix peut chuter à 0, ce qui est l'un des risques associés à l'achat de Bitcoin.}

\item \textbf{Question 192:} B
\\ \textit{Avant d'acheter du Bitcoin, vous devriez avoir un portefeuille sécurisé où vous stockerez vos clés.}

\item \textbf{Question 193:} C
\\ \textit{La technique du Dollar Cost Average implique d'acheter de petites quantités de bitcoin à intervalles réguliers. Cette méthode permet de lisser les fluctuations de prix au fil du temps et assure une croissance continue de la quantité de bitcoin possédée.}

\item \textbf{Question 194:} D
\\ \textit{Un achat spontané est utilisé pour obtenir immédiatement une exposition au Bitcoin. Par exemple, cela pourrait signifier acheter pendant un krach ou profiter d'un bonus.}

\item \textbf{Question 195:} B
\\ \textit{Les réglementations Connaître Votre Client (KYC) exigent que les utilisateurs fournissent une identification pour combattre le financement du terrorisme, l'évasion fiscale et le blanchiment d'argent.}

\item \textbf{Question 196:} D
\\ \textit{Il existe plusieurs manières d'acquérir des bitcoins sans KYC, dont l'une est les distributeurs automatiques de Bitcoin.}

\item \textbf{Question 197:} A
\\ \textit{Après avoir acheté des bitcoins, on doit immédiatement retirer les UTXOs des plateformes d'échange pour minimiser les risques de piratage et de blocage des fonds.}

\item \textbf{Question 198:} D
\\ \textit{Lors du choix d'une plateforme DCA, le bon choix dépendra davantage de sa disponibilité dans votre pays.}

\item \textbf{Question 199:} C
\\ \textit{Consolider vos UTXOs dans vos portefeuilles de temps en temps est essentiel pour gérer efficacement vos bitcoins et éviter des frais inutiles lors des transactions.}

\item \textbf{Question 200:} D
\\ \textit{Le FOMO (Fear of Missing Out) et le FUD (Fear, Uncertainty, Doubt) sont des émotions qui peuvent conduire à des prises de décision impulsives et potentiellement nocives lors de l'achat de Bitcoin.}

\item \textbf{Question 201:} D
\\ \textit{Laisser vos bitcoins sur une plateforme d'échange peut vous exposer aux risques de piratage et de blocage des fonds.}

\item \textbf{Question 202:} A
\\ \textit{Ne pas suivre les réglementations de votre pays lors de l'achat de Bitcoin peut vous mettre en danger.}

\item \textbf{Question 203:} D
\\ \textit{L'adoption du Bitcoin, comme toute nouvelle technologie, suit une courbe en S. C'est un schéma courant pour l'adoption de nouvelles technologies, où l'adoption commence lentement, puis s'accélère avant de ralentir à mesure que la saturation est atteinte.}

\item \textbf{Question 204:} B
\\ \textit{El Salvador a pris la décision audacieuse d'adopter intégralement le Bitcoin, en le déclarant monnaie légale.}

\item \textbf{Question 205:} C
\\ \textit{L'essor de Bitcoin oblige les entreprises, les universités, les régulateurs et les individus à prendre en compte cette nouvelle technologie. De nouveaux outils doivent être créés, les services doivent être adaptés et l'innovation doit continuer pour assurer leur survie.}

\item \textbf{Question 206:} D
\\ \textit{Bitcoin est décrit comme un sujet multidisciplinaire car il soulève des questions liées à divers domaines, incluant la cryptographie, la théorie des jeux, l'économie et la politique monétaire, l'informatique, la philosophie, l'énergie, les lois et la régulation.}

\item \textbf{Question 207:} D
\\ \textit{Selon ce cours, l'objectif principal de la révolution Bitcoin est de créer un monde où la souveraineté humaine est un droit, la vie privée est respectée par défaut, et l'argent n'est pas manipulé.}

\item \textbf{Question 208:} B
\\ \textit{Milton Friedman, un économiste renommé, a prédit en 1999 qu'un e-cash fiable serait bientôt développé, une méthode permettant sur internet de transférer des fonds de A à B, sans que A connaisse B ou vice-versa.}

\item \textbf{Question 209:} A
\\ \textit{Une tendance générale à la hausse pour Bitcoin fait référence à une augmentation constante de son prix sur une certaine période. Il est important de noter que, bien qu'une tendance à la hausse suggère une trajectoire positive, cela ne garantit pas que le prix continuera de monter sans aucune fluctuation.}

\item \textbf{Question 210:} D
\\ \textit{Certains pays ont interdit et criminalisé l'utilisation du Bitcoin, ajoutant à la complexité de son adoption en fonction des cultures, des époques, et des nations. Cependant, le Salvador a été assez courageux pour l'adopter comme monnaie légale.}

\item \textbf{Question 211:} A
\\ \textit{Le Bitcoin est un sujet multidisciplinaire qui soulève des questions liées à divers domaines, incluant la cryptographie, la théorie des jeux, l'économie et la politique monétaire, l'informatique, la philosophie, l'énergie, les lois et la réglementation.}

\item \textbf{Question 212:} A
\\ \textit{Selon ce cours, l'adoption de Bitcoin est virale, suggérant qu'elle se propage rapidement et largement d'une personne à une autre, un peu comme un virus, mais dans le bon sens.}

\item \textbf{Question 213:} C
\\ \textit{Selon ce cours, le Bitcoin ne va pas disparaître. Au contraire, la révolution ne fait que commencer.}

\item \textbf{Question 214:} C
\\ \textit{Le cours suggère qu'il y a tant à explorer avec Bitcoin qu'il est compliqué d'assimiler tout d'un coup, donc on devrait prendre son temps.}

\item \textbf{Question 215:} C
\\ \textit{Le Lightning Network est un réseau de paiement qui utilise le protocole Bitcoin pour permettre des transactions rapides. Ce n'est pas une cryptomonnaie séparée, ni une méthode pour miner du Bitcoin.}

\item \textbf{Question 216:} D
\\ \textit{Le problème de scalabilité fait référence au défi de mettre en place un système monétaire capable de fournir un nombre toujours croissant de transactions par seconde à mesure qu'il est adopté.}

\item \textbf{Question 217:} A
\\ \textit{Le trilemme de la blockchain fait référence au défi selon lequel un protocole basé sur une blockchain ne peut satisfaire que deux critères parmi la décentralisation, la sécurité et la scalabilité.}

\item \textbf{Question 218:} D
\\ \textit{Le Lightning Network permet des transactions Bitcoin instantanées et à faible coût, résolvant ainsi le problème de scalabilité de Bitcoin.}

\item \textbf{Question 219:} D
\\ \textit{Le Lightning Network a été validé et implémenté pour la première fois en 2017.}

\item \textbf{Question 220:} B
\\ \textit{Le Lightning Network agit comme un réseau de canaux de paiement, permettant des transactions instantanées avec des frais réduits pour l'expéditeur.}

\item \textbf{Question 221:} D
\\ \textit{Les services traditionnels de transfert d'argent tels que Western Union, les banques centrales, Visa et Mastercard pourraient disparaître s'ils n'adoptent pas la technologie Lightning Network.}

\item \textbf{Question 222:} B
\\ \textit{Les frais pour les transactions sur le Lightning Network sont extrêmement bas, souvent inférieurs à 0,5 \%.}

\item \textbf{Question 223:} B
\\ \textit{Les transactions sur le Lightning Network sont sécurisées par la cryptographie et indirectement par l'énergie consommée par les mineurs sur Bitcoin.}

\item \textbf{Question 224:} B
\\ \textit{Le Lightning Network permet de réaliser des transactions entre individus qui n'ont pas de connexion directe de canal.}

\item \textbf{Question 225:} C
\\ \textit{L'intervalle de temps moyen entre la création de deux blocs dans le protocole Bitcoin est de 10 minutes.}

\item \textbf{Question 226:} B
\\ \textit{La taille d'un bloc dans le protocole Bitcoin est de 4MB.}

\item \textbf{Question 227:} D
\\ \textit{Le Lightning Network est une solution de deuxième couche au protocole Bitcoin qui vise à faciliter les micro-transactions, qui sont des transactions de très faible valeur qui seraient autrement impraticables en raison des frais plus élevés et des longs temps de confirmation sur la blockchain Bitcoin.}

\item \textbf{Question 228:} D
\\ \textit{Le Lightning Network ouvre la porte à une large gamme d'applications potentielles pour Bitcoin qui étaient auparavant hors de portée en raison des contraintes nécessaires pour assurer la sécurité et la décentralisation de Bitcoin. Parmi ces applications dans la vie quotidienne, nous pouvons mentionner la facturation instantanée dans le commerce.}

\item \textbf{Question 229:} D
\\ \textit{Le Lightning Network utilise des satoshis (également appelés sats), le décimal de bitcoin, pour fonctionner.}

\item \textbf{Question 230:} D
\\ \textit{Le concept de skin in the game est une idée qui a récemment gagné en popularité dans l'industrie du jeu vidéo. Il fait essentiellement référence à l'idée que les développeurs et les parties prenantes ont un investissement personnel ou un enjeu dans le succès d'un jeu, ce qui aligne leurs intérêts avec la livraison d'un produit de qualité.}

\item \textbf{Question 231:} C
\\ \textit{Le Lightning Network permet aux joueurs de miser de très petites sommes d'argent lorsqu'ils jouent à des jeux, comme quelques satoshis. Cela permet d'établir un enjeu qui stimule la compétition tout en augmentant significativement le coût de déploiement de bots.}

\item \textbf{Question 232:} B
\\ \textit{Avec le Lightning Network, nous pouvons effectuer des micro-transactions chaque minute, ce qui ouvre la porte à l'expérimentation de modèles économiques où les consommateurs paient pour le contenu en fonction de ce qu'ils consomment.}

\item \textbf{Question 233:} A
\\ \textit{Avec les micro-transactions sur le Lightning Network, les utilisateurs peuvent bénéficier de frais de transaction faibles, ce qui rend économique l'envoi de très petites quantités de Bitcoin. Ils peuvent également avoir des paiements instantanés, permettant des transactions en temps réel. Enfin, ils peuvent obtenir une scalabilité, permettant de gérer un grand nombre de transactions simultanément. Ainsi, la mauvaise réponse est les temps de paiement longs.}

\item \textbf{Question 234:} B
\\ \textit{Avec le Lightning Network, les créateurs de contenu seraient incités à produire du contenu de bonne qualité pour retenir l'attention des utilisateurs. Les utilisateurs, à leur tour, ne paieraient que pour le contenu qu'ils consomment, éliminant ainsi les frais d'abonnement préalables.}

\item \textbf{Question 235:} B
\\ \textit{Une plateforme DCA fait référence à une plateforme qui facilite le Dollar-Cost Averaging (DCA), une stratégie d'investissement où un individu investit un montant fixe d'argent dans un actif particulier à intervalles réguliers, indépendamment du prix de l'actif.}

\item \textbf{Question 236:} C
\\ \textit{Avec le réseau Lightning, il est concevable d'utiliser ce système pour la location de biens. Les utilisateurs pourraient être facturés à la minute, ou même à la seconde, pour le temps qu'ils utilisent les biens.}

\item \textbf{Question 237:} C
\\ \textit{Le Lightning Network ouvre la porte à une multitude de cas d'utilisation passionnants pour les utilisateurs de Bitcoin. Les modèles économiques et les opportunités commerciales qui en résultent sont nombreux et variés, incluant de nouveaux modèles économiques où les consommateurs paient pour le contenu en fonction de ce qu'ils consomment.}

\item \textbf{Question 238:} B
\\ \textit{Une manière de tester le Lightning Network par vous-même est d'essayer l'application de podcast Fountain, qui vous permet d'être récompensé avec quelques sats pour l'écoute de vos podcasts préférés.}

\item \textbf{Question 239:} A
\\ \textit{Une implication potentielle de l'avancement des technologies IA est que plus de 80 \% des emplois disparaîtront en raison de l'IA et de l'automatisation.}

\item \textbf{Question 240:} D
\\ \textit{Bitcoin, à l'instar d'internet pour les moyens de communication, est une révolution technologique qui crée la possibilité de nouveaux modes d'organisation à grande échelle, nous apportant la possibilité d'échanger de la valeur sans aucun tiers de confiance.}

\item \textbf{Question 241:} B
\\ \textit{Le cours soulève plusieurs questions sur l'avenir de la finance, dont l'une est qui devrait détenir, autoriser et tracer l'argent que nous utilisons.}

\item \textbf{Question 242:} C
\\ \textit{Le Bitcoin est pseudonyme et peut être échangé partout dans le monde sans l'autorisation d'aucune entité.}

\item \textbf{Question 243:} C
\\ \textit{L'acceptation de ces nouvelles technologies pourrait générer d'énormes économies à l'échelle mondiale.}

\item \textbf{Question 244:} B
\\ \textit{La question de savoir qui devrait contrôler le système bancaire est cruciale parce que les règles du jeu bancaire ne sont pas transparentes et compréhensibles pour tous.}

\item \textbf{Question 245:} C
\\ \textit{Tolérer la censure peut détruire la liberté d'expression et le droit de réunion.}

\item \textbf{Question 246:} D
\\ \textit{Il y a 2,4 milliards de personnes dans le monde sans compte bancaire, et le Bitcoin permet l'égalité dans les transactions, indépendamment de leur statut social ou de leur position politique.}

\item \textbf{Question 247:} A
\\ \textit{Le protocole Bitcoin est neutre envers tous ses utilisateurs.}

\item \textbf{Question 248:} D
\\ \textit{Les utilisateurs peuvent obtenir des bitcoins en les acceptant en échange ou en les achetant via des plateformes régulées ou non régulées.}

\item \textbf{Question 249:} C
\\ \textit{Satoshi Nakamoto a créé Bitcoin en 2008 pour proposer le changement du système financier.}

\item \textbf{Question 250:} C
\\ \textit{Le Bitcoin permet l'émancipation du système bancaire.}


\end{enumerate}

\end{document}
